% =====================================================================
%  Chapitre 2 : suite de la section 1 (mécanique quantique)
%  Fichier importé par main.tex avec % =====================================================================
%  Chapitre 2 : suite de la section 1 (mécanique quantique)
%  Fichier importé par main.tex avec % =====================================================================
%  Chapitre 2 : suite de la section 1 (mécanique quantique)
%  Fichier importé par main.tex avec % =====================================================================
%  Chapitre 2 : suite de la section 1 (mécanique quantique)
%  Fichier importé par main.tex avec \include{chapitre2}.
%  Il PROLONGE chapitre1.tex : pas de \section ici, les sous-sections
%  1.6 à 1.9 font suite à 1.1 à 1.5 (la suite est dans chapitre3.tex).
%  Pas de préambule non plus : les
%  environnements (definition, resultat, remarque), les macros (\ii, \ee,
%  \dd, \pd, \pdd, \Hop, \ket, \bra, \braket) et les couleurs (insarouge,
%  insagris) sont définis dans main.tex.
%  Convention de signe : p = -i hbar d/dx, cohérente avec l'onde plane
%  e^{i(kx-wt)} de la section 1.3 (les notes de cours donnaient +i hbar).
% =====================================================================

% Macros de Dirac (déjà définies dans main.tex récent ; ce bloc évite l'erreur
% « Undefined control sequence \ket » si main.tex est une ancienne version).
\providecommand{\ket}[1]{|#1\rangle}
\providecommand{\bra}[1]{\langle #1|}
\providecommand{\braket}[2]{\langle #1|#2\rangle}

% ---------------------------------------------------------------------
\subsection{Notation de Dirac}\label{sec:dirac}
% ---------------------------------------------------------------------

Un état s'écrit comme un vecteur colonne, par exemple
$\begin{pmatrix} a_1\\ a_2\end{pmatrix}\in\mathbb{C}^2$ : c'est un élément d'un
espace de Hilbert de dimension~2, noté $\mathcal{H}^2$. Dirac le note simplement
$\ket{a}$ et évite d'écrire les composantes à chaque ligne.

\begin{definition}{Notation de Dirac : ket, bra et bra-ket}
\begin{center}
\begin{tabular}{@{}r@{\hspace{0.5em}}l@{\hspace{2em}}l@{}}
  {\Large\color{insarouge}$\ket{a}$}
    & $=\begin{pmatrix}a_1\\ a_2\end{pmatrix}\in\mathcal{H}^2$
    & \textbf{ket} : vecteur colonne \\[1.6ex]
  {\Large\color{insarouge}$\bra{a}$}
    & $=\big(\ket{a}\big)^\dagger=\begin{pmatrix}a_1^* & a_2^*\end{pmatrix}$
    & \textbf{bra} : vecteur ligne \\[1.6ex]
  {\Large\color{insarouge}$\braket{a}{b}$}
    & $=\begin{pmatrix}a_1^* & a_2^*\end{pmatrix}\begin{pmatrix}b_1\\ b_2\end{pmatrix}
       =a_1^*b_1+a_2^*b_2$
    & \textbf{bra-ket} : produit scalaire \\
\end{tabular}
\end{center}
Le bra est le transposé \emph{conjugué} du ket. Le nom vient de l'anglais
\emph{bracket} (crochet) : en accolant $\bra{a}$~(\emph{bra}) et $\ket{b}$~(\emph{ket})
on obtient $\braket{a}{b}$~(\emph{bracket}), un nombre complexe.

Propriétés : $\braket{b}{a}=\braket{a}{b}^*$ et
$\braket{a}{a}=|a_1|^2+|a_2|^2\in\mathbb{R}_+$ (norme au carré).
\end{definition}

Dans l'autre ordre, le produit d'un ket par un bra est une matrice (un opérateur) :
\[
  \ket{a}\bra{b}
  =\begin{pmatrix}a_1\\ a_2\end{pmatrix}\begin{pmatrix}b_1^* & b_2^*\end{pmatrix}
  =\begin{pmatrix}a_1b_1^* & a_1b_2^*\\ a_2b_1^* & a_2b_2^*\end{pmatrix}.
\]

\begin{remarque}[Pourquoi le conjugué ?]
Pour $\ket{a}=\dfrac{1}{\sqrt2}\begin{pmatrix}1\\ \ii\end{pmatrix}$, on a
$\bra{a}=\dfrac{1}{\sqrt2}\begin{pmatrix}1 & -\ii\end{pmatrix}$ et
$\braket{a}{a}=\tfrac12\,(1+1)=1$. Sans conjugaison on trouverait
$\tfrac12\,(1+\ii^2)=0$, ce qui n'est pas une norme.
\end{remarque}

Un état physique est \emph{normalisé} :
\[
  \braket{\Psi}{\Psi}=1 .
\]
Pour une particule dans l'espace, $\ket{\Psi}$ est une fonction de $x$ (dimension
infinie) et le produit scalaire devient une intégrale,
$\braket{\varphi}{\psi}=\int\varphi^*\psi\,\dd x$. On retrouve la condition de
normalisation de Born (section~\ref{sec:born}) :
\[
  \braket{\Psi}{\Psi} = \int_{-\infty}^{+\infty}\Psi^*\Psi\,\dd x = 1 .
\]
L'espace de Hilbert et le formalisme complet sont repris plus loin dans le cours.

% ---------------------------------------------------------------------
\subsection{Opérateurs quantiques}\label{sec:operateurs}
% ---------------------------------------------------------------------

À chaque grandeur physique (position, quantité de mouvement, énergie) on associe
un \emph{opérateur} $\hat O$ : une règle qui transforme une fonction d'onde en une
autre, $\Psi\mapsto\hat O\Psi$. Le point de départ est l'onde plane de la
section~\ref{sec:construction} : $\Psi=\ee^{\ii(kx-\omega t)}$ vérifie
$-\ii\hbar\,\partial_x\Psi=\hbar k\,\Psi=p\,\Psi$.

\begin{definition}{Opérateur quantité de mouvement}
\[
  \hat p = -\ii\hbar\,\frac{\partial}{\partial x},
  \qquad
  \hat p\,\Psi = -\ii\hbar\,\frac{\partial\Psi}{\partial x}.
\]
L'onde plane est fonction propre de $\hat p$, de valeur propre $p=\hbar k$.
\end{definition}

Les autres opérateurs s'obtiennent en remplaçant $p$ par $\hat p$ dans
l'expression classique. Avec $E_c=\tfrac12 mv^2=p^2/2m$ :
\[
  \hat E_c = \frac{\hat p^{\,2}}{2m}
  = \frac{1}{2m}\Big(-\ii\hbar\frac{\partial}{\partial x}\Big)
                \Big(-\ii\hbar\frac{\partial}{\partial x}\Big)
  = \frac{(-\ii\hbar)^2}{2m}\,\frac{\partial^2}{\partial x^2}
  = -\frac{\hbar^2}{2m}\,\frac{\partial^2}{\partial x^2},
  \qquad
  \hat v = \frac{\hat p}{m} = -\frac{\ii\hbar}{m}\,\frac{\partial}{\partial x}.
\]

\begin{center}
\renewcommand{\arraystretch}{1.6}
\begin{tabular}{@{}lll@{}}
\toprule
\textbf{Grandeur} & \textbf{Classique} & \textbf{Opérateur} \\
\midrule
Position & $x$ & $\hat x = x\times$ \\
Quantité de mouvement & $p$ & $\hat p = -\ii\hbar\,\partial_x$ \\
Vitesse & $v=\dfrac{p}{m}$ & $\hat v = -\dfrac{\ii\hbar}{m}\,\partial_x$ \\
Énergie cinétique & $E_c=\dfrac{p^2}{2m}$ & $\hat E_c = -\dfrac{\hbar^2}{2m}\,\partial_x^2$ \\
Énergie potentielle & $V(x)$ & $\hat V = V(x)\times$ \\
Énergie totale & $E=E_c+V$ & $\Hop = \hat E_c+\hat V = -\dfrac{\hbar^2}{2m}\,\partial_x^2+V(x)$ \\
\bottomrule
\end{tabular}
\end{center}

L'opérateur hamiltonien $\Hop$ représente l'énergie totale. Avec $E\to\ii\hbar\,\partial_t$,
l'équation de Schrödinger de la section~\ref{sec:construction} n'est que la relation
$E=E_c+V$ écrite avec des opérateurs : $\ii\hbar\,\partial_t\Psi=\Hop\Psi$.

\begin{remarque}[Convention de signe]
Le signe de $\hat p$ dépend du choix de l'onde plane $\ee^{\ii(kx-\omega t)}$. La
convention conjuguée $\ee^{-\ii(kx-\omega t)}$ inverse les deux signes
($\hat p=+\ii\hbar\,\partial_x$ et $E\to-\ii\hbar\,\partial_t$) sans changer les
résultats physiques. Comme $\hat p$ intervient au carré, $\hat E_c$ est le même dans
les deux conventions.
\end{remarque}

\begin{definition}{Valeur moyenne d'un opérateur}
Pour une fonction d'onde normalisée $\Psi$,
\[
  \langle\hat O\rangle
  = \int_{-\infty}^{+\infty}\Psi^*(x,t)\;\hat O\,\Psi(x,t)\,\dd x
  = \bra{\Psi}\hat O\ket{\Psi}.
\]
C'est la moyenne des résultats obtenus en mesurant la grandeur sur un grand nombre
de systèmes identiques, tous préparés dans l'état $\Psi$. L'opérateur agit sur
$\Psi$ seul, à droite de $\Psi^*$. Pour la position,
$\langle\hat x\rangle=\int x\,|\Psi|^2\,\dd x$ : c'est l'espérance de la densité
de probabilité de Born.

L'\emph{incertitude} (écart type) sur la grandeur associée à $\hat O$ mesure la
dispersion des résultats autour de cette moyenne :
\[
  \sigma_O=\sqrt{\langle\hat O^{\,2}\rangle-\langle\hat O\rangle^2}\,,
  \qquad
  \langle\hat O^{\,2}\rangle=\int\Psi^*\,\hat O^{\,2}\,\Psi\,\dd x .
\]
Si $\Psi$ est une fonction propre de $\hat O$ ($\hat O\Psi=\lambda\Psi$), alors
$\langle\hat O\rangle=\lambda$, $\langle\hat O^{\,2}\rangle=\lambda^2$ et
$\sigma_O=0$ : la grandeur a une valeur certaine.
\end{definition}

% ---------------------------------------------------------------------
\subsection{Solutions de l'équation de Schrödinger}\label{sec:solutions}
% ---------------------------------------------------------------------

\subsubsection{Du potentiel à la fonction d'onde}

En mécanique classique, l'énergie potentielle $V(x,t)$ fixe la force
$F=-\partial V/\partial x$, puis $\dd p/\dd t=F$ donne la trajectoire. En mécanique
quantique, c'est l'équation de Schrödinger qui joue ce rôle : connaissant $V$, on la
résout pour obtenir $\Psi(x,t)$.

\begin{resultat}{ESDT : équation de Schrödinger dépendante du temps}
\[
  \ii\hbar\,\pd{\Psi(x,t)}{t} = \Hop\,\Psi(x,t),
  \qquad
  \Hop = -\frac{\hbar^2}{2m}\,\pdd{}{x} + V(x,t).
\]
\end{resultat}

Si l'énergie potentielle ne dépend pas du temps, $V(x,t)=V(x)$, la force
$F=-\dd V/\dd x$ (dérivée ordinaire) est constante dans le temps et l'énergie se
conserve. L'équation se résout alors par séparation des variables.

\subsubsection{Séparation des variables}

On cherche des solutions produit d'une fonction de $x$ seul et d'une fonction de $t$
seul :
\[
  \Psi(x,t)=\psi(x)\,\varphi(t).
\]
Dans l'ESDT, les dérivées partielles deviennent des dérivées ordinaires :
\[
  \ii\hbar\,\psi(x)\,\frac{\dd\varphi}{\dd t}
  = \varphi(t)\left[-\frac{\hbar^2}{2m}\,\frac{\dd^2\psi}{\dd x^2}+V(x)\,\psi(x)\right].
\]
On divise par $\psi(x)\,\varphi(t)$ :
\[
  \underbrace{\ii\hbar\,\frac{1}{\varphi}\,\frac{\dd\varphi}{\dd t}}_{\text{ne dépend que de }t}
  \;=\;
  \underbrace{\frac{1}{\psi}\left[-\frac{\hbar^2}{2m}\,\frac{\dd^2\psi}{\dd x^2}+V(x)\,\psi\right]}_{\text{ne dépend que de }x}.
\]
Les deux membres dépendent de variables indépendantes et sont pourtant égaux pour
tout $(x,t)$ : ils sont égaux à une même constante, notée $E$. On obtient deux
équations.

\paragraph{Partie temporelle.} $\ii\hbar\,\dd\varphi/\dd t=E\,\varphi$ donne
\[
  \varphi(t)=\ee^{-\ii Et/\hbar}=\ee^{-\ii\omega t},
  \qquad \omega=\frac{E}{\hbar},
\]
la constante multiplicative étant absorbée dans $\psi$. On retrouve $E=\hbar\omega$ :
la constante $E$ est bien l'énergie de la particule.

\paragraph{Partie spatiale.} C'est l'ESIT ci-dessous.

\begin{resultat}{ESIT : équation de Schrödinger indépendante du temps}
\[
  -\frac{\hbar^2}{2m}\,\frac{\dd^2\psi(x)}{\dd x^2}+V(x)\,\psi(x)=E\,\psi(x)
  \qquad\Longleftrightarrow\qquad
  \Hop\,\psi(x)=E\,\psi(x).
\]
C'est un \emph{problème aux valeurs propres} : $E$ est une valeur propre de $\Hop$
(énergie permise) et $\psi$ la fonction propre associée. Chaque solution
$(E_n,\psi_n)$ donne un \emph{état stationnaire}
\[
  \Psi_n(x,t)=\psi_n(x)\,\ee^{-\ii E_nt/\hbar}.
\]
\end{resultat}

Dans un état stationnaire, $|\Psi_n(x,t)|^2=|\psi_n(x)|^2$ ne dépend pas du temps :
le facteur temporel est une phase globale, non observable. Les conditions physiques
($\psi$ continue et normalisable) ne sont satisfaites que pour certaines valeurs de
$E$ : c'est l'origine de la quantification de l'énergie, illustrée ci-dessous.

Comme l'ESDT est linéaire (hypothèse~(c) de la section~\ref{sec:construction}), toute
superposition $\Psi=\sum_n c_n\,\Psi_n$ est aussi solution, normalisée si
$\sum_n|c_n|^2=1$ et si les $\psi_n$ sont orthonormées. Une telle superposition
n'est en général pas stationnaire : $|\Psi|^2$ contient des termes croisés en
$\ee^{\ii(E_m-E_n)t/\hbar}$.

\subsubsection{Exemple : puits de potentiel infini}

Une particule de masse $m$ (un électron par exemple) est enfermée entre deux parois
impénétrables, en $x=0$ et $x=a$ :
\[
  V(x)=
  \begin{cases}
    0 & \text{si } 0<x<a,\\
    +\infty & \text{ailleurs.}
  \end{cases}
\]
La particule ne peut pas se trouver hors du puits : $\psi=0$ pour $x\le0$ et
$x\ge a$. La fonction $\psi$ est continue (mais pas sa dérivée, à cause du saut
infini de $V$), d'où les conditions aux limites
\[
  \psi(0)=\psi(a)=0 .
\]

\paragraph{Opérateur hamiltonien.} Dans le puits $V=0$, donc $\Hop$ se réduit à
l'opérateur énergie cinétique :
\[
  \Hop=\hat E_c=-\frac{\hbar^2}{2m}\,\frac{\dd^2}{\dd x^2}\qquad(0<x<a).
\]

\paragraph{Résolution de l'ESIT.} Avec $V=0$, l'ESIT s'écrit
$\psi''=-k^2\psi$ avec $k=\sqrt{2mE}/\hbar$. On prend $E>0$ : pour $E\le0$, la seule
solution vérifiant $\psi(0)=\psi(a)=0$ est $\psi\equiv0$. La solution générale est
\[
  \psi(x)=A\sin(kx)+B\cos(kx).
\]
\begin{itemize}[nosep]
  \item $\psi(0)=0$ donne $B=0$ ;
  \item $\psi(a)=A\sin(ka)=0$ avec $A\neq0$ (sinon $\psi\equiv0$ : pas de
        particule) donne $ka=n\pi$, soit
        $k_n=\dfrac{n\pi}{a}$ avec $n=1,2,3,\dots$
        ($n=0$ donne $\psi\equiv0$, et $n<0$ redonne les mêmes états au signe près).
\end{itemize}
Cela revient à $\lambda_n=2\pi/k_n=2a/n$, soit $a=n\,\lambda_n/2$ : $n$ demi-longueurs
d'onde tiennent dans le puits, comme pour une onde stationnaire.

\begin{resultat}{Énergie quantifiée}
\[
  E_n=\frac{\hbar^2k_n^2}{2m}=\frac{n^2\pi^2\hbar^2}{2ma^2}=\frac{n^2h^2}{8ma^2}=n^2E_1,
  \qquad n=1,2,3,\dots
\]
\textbf{Énergie fondamentale} (niveau $n=1$, le plus bas) :
$\displaystyle E_1=\frac{\pi^2\hbar^2}{2ma^2}=\frac{h^2}{8ma^2}>0$.
\end{resultat}

\begin{remarque}[Lecture du résultat]
\begin{itemize}[nosep]
  \item Les énergies sont \emph{discrètes} : $E_{n+1}-E_n=(2n+1)E_1$. Comme
        $E_n\propto1/a^2$, plus le puits est étroit, plus les énergies sont grandes.
  \item $E_1\neq0$ : une particule confinée ne peut pas être au repos (énergie de
        point zéro), conséquence du principe d'incertitude de Heisenberg.
  \item Ordre de grandeur : pour un électron avec $a=1$~nm,
        $E_1\approx0{,}38$~eV et $E_2=4E_1\approx1{,}5$~eV. Pour $a=1$~cm,
        $E_1\approx4\times10^{-15}$~eV : la quantification est indétectable.
\end{itemize}
\end{remarque}

\paragraph{Normalisation.} Le module de $\ee^{-\ii E_nt/\hbar}$ vaut 1, donc
$|\Psi_n|^2=A^2\sin^2(n\pi x/a)$ et la condition de Born impose
\[
  \int_0^a|\Psi_n|^2\,\dd x
  = A^2\int_0^a\frac{1-\cos\!\big(2n\pi x/a\big)}{2}\,\dd x
  = A^2\,\frac{a}{2}=1
  \quad\Longrightarrow\quad A=\sqrt{\frac{2}{a}},
\]
car l'intégrale du cosinus sur $n$ périodes est nulle.

\begin{resultat}{États stationnaires du puits infini}
\[
  \Psi_n(x,t)=\sqrt{\frac{2}{a}}\;\sin\!\Big(\frac{n\pi x}{a}\Big)\,\ee^{-\ii E_nt/\hbar}
  \quad(0<x<a),
  \qquad
  |\Psi_n|^2=\frac{2}{a}\,\sin^2\!\Big(\frac{n\pi x}{a}\Big).
\]
$\psi_n$ possède $n-1$ nœuds à l'intérieur du puits. Pour $n=1$ (état fondamental,
d'énergie $E_1$) :
\[
  \Psi_1(x,t)=\sqrt{\frac{2}{a}}\;\sin\!\Big(\frac{\pi x}{a}\Big)\,\ee^{-\ii E_1t/\hbar}.
\]
\end{resultat}

\begin{figure}[!ht]
\centering
\begin{tikzpicture}[>=Stealth, x=1cm, y=0.95cm]
  % ---------- (a) niveaux d'énergie et fonctions propres ----------
  \begin{scope}
    \node[above, font=\small] at (2,5.0) {(a) niveaux $E_n$ et fonctions propres $\psi_n$};
    \draw[->] (-2.5,0) -- (-2.5,4.8) node[above] {$E$};
    \draw[very thick, insarouge] (0,4.4) -- (0,0) -- (4,0) -- (4,4.4);
    \node[insarouge, right, font=\small] at (4.05,3.9) {$V=+\infty$};
    \node[insarouge, below, font=\small] at (2,0) {$V=0$};
    \node[below, font=\small] at (0,-0.35) {$0$};
    \node[below, font=\small] at (4,-0.35) {$a$};
    \node[left, font=\small] at (-0.1,0.4) {$E_1$};
    \node[left, font=\small] at (-0.1,1.6) {$E_2=4E_1$};
    \node[left, font=\small] at (-0.1,3.6) {$E_3=9E_1$};
    \foreach \n/\lev in {1/0.4, 2/1.6, 3/3.6} {
      \draw[dashed, insagris!70] (0,\lev) -- (4,\lev);
      \draw[thick, insagris, domain=0:1, samples=80]
        plot ({4*\x}, {\lev+0.5*sin(\n*180*\x)});
    }
  \end{scope}
  % ---------- (b) densités de probabilité ----------
  \begin{scope}[xshift=7cm]
    \node[above, font=\small] at (2,5.0) {(b) densités de probabilité $|\psi_n|^2$};
    \draw[->] (-0.3,0) -- (4.6,0) node[right] {$x$};
    \draw[thick, insarouge] (0,0) -- (0,4.4) (4,0) -- (4,4.4);
    \node[below, font=\small] at (0,0) {$0$};
    \node[below, font=\small] at (4,0) {$a$};
    \node[below, font=\small] at (2,0) {$a/2$};
    \foreach \n/\lev in {1/0.4, 2/1.6, 3/3.6} {
      \draw[dashed, insagris!70] (0,\lev) -- (4,\lev);
      \fill[insarouge!20] (0,\lev) -- plot[domain=0:1, samples=80]
        ({4*\x}, {\lev+0.6*(sin(\n*180*\x))^2}) -- (4,\lev) -- cycle;
      \draw[thick, insarouge, domain=0:1, samples=80]
        plot ({4*\x}, {\lev+0.6*(sin(\n*180*\x))^2});
    }
    \draw[dotted, thick, insagris] (2,0) -- (2,4.4);
    \node[above, font=\small] at (2,4.4) {$\langle\hat x\rangle=a/2$};
  \end{scope}
\end{tikzpicture}
\caption{Puits de potentiel infini de largeur $a$. (a) Niveaux d'énergie
$E_n=n^2E_1$ et fonctions propres $\psi_n$. (b) Densités de probabilité
$|\psi_n|^2$, symétriques par rapport à $a/2$.}
\end{figure}

\subsubsection{Position moyenne dans le puits}

D'après la définition de la valeur moyenne, avec $\hat x=x\times$ :
\[
  \langle\hat x\rangle
  = \int_0^a\Psi_n^*(x,t)\;x\;\Psi_n(x,t)\,\dd x
  = \frac{2}{a}\int_0^a x\,\sin^2\!\Big(\frac{n\pi x}{a}\Big)\dd x
  = \frac{1}{a}\int_0^a x\Big(1-\cos\frac{2n\pi x}{a}\Big)\dd x .
\]
Le premier terme vaut $\dfrac1a\cdot\dfrac{a^2}{2}=\dfrac a2$. Pour le second, une
intégration par parties donne
\[
  \int_0^a x\cos\frac{2n\pi x}{a}\,\dd x
  =\Big[\frac{a\,x}{2n\pi}\sin\frac{2n\pi x}{a}\Big]_0^a
   +\frac{a^2}{4n^2\pi^2}\Big[\cos\frac{2n\pi x}{a}\Big]_0^a = 0,
\]
puisque $\sin(2n\pi)=0$ et $\cos(2n\pi)=1$.

\begin{resultat}{Position moyenne}
\[
  \langle\hat x\rangle=\frac{a}{2}\qquad\text{pour tout }n\text{ et à tout instant.}
\]
\end{resultat}

Les phases $\ee^{\pm\ii E_nt/\hbar}$ se compensent dans $\Psi^*\Psi$. Le résultat se
voit d'ailleurs sans calcul, car $|\Psi_n|^2$ est symétrique par rapport à $a/2$.

\begin{remarque}[Moyenne et position la plus probable]
Pour $n=2$, la densité de probabilité s'annule en $x=a/2$ alors que
$\langle\hat x\rangle=a/2$ : la valeur moyenne n'est pas la position la plus
probable, et la particule n'est jamais trouvée là.
\end{remarque}

\paragraph{Incertitude sur l'énergie.} Comme $\Hop\Psi_n=E_n\Psi_n$, on a aussi
$\Hop^{2}\Psi_n=E_n^2\Psi_n$, donc $\langle\Hop\rangle=E_n$, $\langle\Hop^{2}\rangle=E_n^2$
et
\[
  \sigma_E=\sqrt{\langle\Hop^{2}\rangle-\langle\Hop\rangle^2}=\sqrt{E_n^2-E_n^2}=0 .
\]
Dans un état propre (stationnaire) l'énergie est parfaitement déterminée, alors que la
position ne l'est pas. Que se passe-t-il si la particule n'est pas dans un état propre ?
C'est l'objet de la section~\ref{sec:superposition}.

% ---------------------------------------------------------------------
\subsection{Bilan : opérateurs et puits infini}
% ---------------------------------------------------------------------

\begin{center}
\small\renewcommand{\arraystretch}{1.25}
\begin{tabular}{@{}ll@{}}
\toprule
\textbf{Notion} & \textbf{Formule} \\
\midrule
Bra-ket & $\braket{a}{b}=a_1^*b_1+a_2^*b_2$, \quad
          $\braket{\Psi}{\Psi}=\displaystyle\int\Psi^*\Psi\,\dd x=1$ \\
Opérateurs & $\hat p=-\ii\hbar\,\partial_x$, \quad
             $\hat E_c=\dfrac{\hat p^{\,2}}{2m}=-\dfrac{\hbar^2}{2m}\,\partial_x^2$, \quad
             $\Hop=\hat E_c+V(x)$ \\
Valeur moyenne, incertitude & $\langle\hat O\rangle=\displaystyle\int\Psi^*\,\hat O\,\Psi\,\dd x$,
          \quad $\sigma_O=\sqrt{\langle\hat O^{\,2}\rangle-\langle\hat O\rangle^2}$ \\
ESDT & $\ii\hbar\,\partial_t\Psi=\Hop\,\Psi$ \\
Séparation ($V=V(x)$) & $\Psi=\psi(x)\,\ee^{-\ii Et/\hbar}$, \quad ESIT : $\Hop\,\psi=E\,\psi$ \\
Puits infini & $\psi_n=\sqrt{2/a}\,\sin(n\pi x/a)$, \quad
               $E_n=\dfrac{n^2\pi^2\hbar^2}{2ma^2}$, \quad $E_1=\dfrac{\pi^2\hbar^2}{2ma^2}$ \\
Moyennes dans $\Psi_n$ & $\langle\hat x\rangle=a/2$, \quad $\langle\Hop\rangle=E_n$, \quad $\sigma_E=0$ \\
\bottomrule
\end{tabular}
\end{center}
.
%  Il PROLONGE chapitre1.tex : pas de \section ici, les sous-sections
%  1.6 à 1.9 font suite à 1.1 à 1.5 (la suite est dans chapitre3.tex).
%  Pas de préambule non plus : les
%  environnements (definition, resultat, remarque), les macros (\ii, \ee,
%  \dd, \pd, \pdd, \Hop, \ket, \bra, \braket) et les couleurs (insarouge,
%  insagris) sont définis dans main.tex.
%  Convention de signe : p = -i hbar d/dx, cohérente avec l'onde plane
%  e^{i(kx-wt)} de la section 1.3 (les notes de cours donnaient +i hbar).
% =====================================================================

% Macros de Dirac (déjà définies dans main.tex récent ; ce bloc évite l'erreur
% « Undefined control sequence \ket » si main.tex est une ancienne version).
\providecommand{\ket}[1]{|#1\rangle}
\providecommand{\bra}[1]{\langle #1|}
\providecommand{\braket}[2]{\langle #1|#2\rangle}

% ---------------------------------------------------------------------
\subsection{Notation de Dirac}\label{sec:dirac}
% ---------------------------------------------------------------------

Un état s'écrit comme un vecteur colonne, par exemple
$\begin{pmatrix} a_1\\ a_2\end{pmatrix}\in\mathbb{C}^2$ : c'est un élément d'un
espace de Hilbert de dimension~2, noté $\mathcal{H}^2$. Dirac le note simplement
$\ket{a}$ et évite d'écrire les composantes à chaque ligne.

\begin{definition}{Notation de Dirac : ket, bra et bra-ket}
\begin{center}
\begin{tabular}{@{}r@{\hspace{0.5em}}l@{\hspace{2em}}l@{}}
  {\Large\color{insarouge}$\ket{a}$}
    & $=\begin{pmatrix}a_1\\ a_2\end{pmatrix}\in\mathcal{H}^2$
    & \textbf{ket} : vecteur colonne \\[1.6ex]
  {\Large\color{insarouge}$\bra{a}$}
    & $=\big(\ket{a}\big)^\dagger=\begin{pmatrix}a_1^* & a_2^*\end{pmatrix}$
    & \textbf{bra} : vecteur ligne \\[1.6ex]
  {\Large\color{insarouge}$\braket{a}{b}$}
    & $=\begin{pmatrix}a_1^* & a_2^*\end{pmatrix}\begin{pmatrix}b_1\\ b_2\end{pmatrix}
       =a_1^*b_1+a_2^*b_2$
    & \textbf{bra-ket} : produit scalaire \\
\end{tabular}
\end{center}
Le bra est le transposé \emph{conjugué} du ket. Le nom vient de l'anglais
\emph{bracket} (crochet) : en accolant $\bra{a}$~(\emph{bra}) et $\ket{b}$~(\emph{ket})
on obtient $\braket{a}{b}$~(\emph{bracket}), un nombre complexe.

Propriétés : $\braket{b}{a}=\braket{a}{b}^*$ et
$\braket{a}{a}=|a_1|^2+|a_2|^2\in\mathbb{R}_+$ (norme au carré).
\end{definition}

Dans l'autre ordre, le produit d'un ket par un bra est une matrice (un opérateur) :
\[
  \ket{a}\bra{b}
  =\begin{pmatrix}a_1\\ a_2\end{pmatrix}\begin{pmatrix}b_1^* & b_2^*\end{pmatrix}
  =\begin{pmatrix}a_1b_1^* & a_1b_2^*\\ a_2b_1^* & a_2b_2^*\end{pmatrix}.
\]

\begin{remarque}[Pourquoi le conjugué ?]
Pour $\ket{a}=\dfrac{1}{\sqrt2}\begin{pmatrix}1\\ \ii\end{pmatrix}$, on a
$\bra{a}=\dfrac{1}{\sqrt2}\begin{pmatrix}1 & -\ii\end{pmatrix}$ et
$\braket{a}{a}=\tfrac12\,(1+1)=1$. Sans conjugaison on trouverait
$\tfrac12\,(1+\ii^2)=0$, ce qui n'est pas une norme.
\end{remarque}

Un état physique est \emph{normalisé} :
\[
  \braket{\Psi}{\Psi}=1 .
\]
Pour une particule dans l'espace, $\ket{\Psi}$ est une fonction de $x$ (dimension
infinie) et le produit scalaire devient une intégrale,
$\braket{\varphi}{\psi}=\int\varphi^*\psi\,\dd x$. On retrouve la condition de
normalisation de Born (section~\ref{sec:born}) :
\[
  \braket{\Psi}{\Psi} = \int_{-\infty}^{+\infty}\Psi^*\Psi\,\dd x = 1 .
\]
L'espace de Hilbert et le formalisme complet sont repris plus loin dans le cours.

% ---------------------------------------------------------------------
\subsection{Opérateurs quantiques}\label{sec:operateurs}
% ---------------------------------------------------------------------

À chaque grandeur physique (position, quantité de mouvement, énergie) on associe
un \emph{opérateur} $\hat O$ : une règle qui transforme une fonction d'onde en une
autre, $\Psi\mapsto\hat O\Psi$. Le point de départ est l'onde plane de la
section~\ref{sec:construction} : $\Psi=\ee^{\ii(kx-\omega t)}$ vérifie
$-\ii\hbar\,\partial_x\Psi=\hbar k\,\Psi=p\,\Psi$.

\begin{definition}{Opérateur quantité de mouvement}
\[
  \hat p = -\ii\hbar\,\frac{\partial}{\partial x},
  \qquad
  \hat p\,\Psi = -\ii\hbar\,\frac{\partial\Psi}{\partial x}.
\]
L'onde plane est fonction propre de $\hat p$, de valeur propre $p=\hbar k$.
\end{definition}

Les autres opérateurs s'obtiennent en remplaçant $p$ par $\hat p$ dans
l'expression classique. Avec $E_c=\tfrac12 mv^2=p^2/2m$ :
\[
  \hat E_c = \frac{\hat p^{\,2}}{2m}
  = \frac{1}{2m}\Big(-\ii\hbar\frac{\partial}{\partial x}\Big)
                \Big(-\ii\hbar\frac{\partial}{\partial x}\Big)
  = \frac{(-\ii\hbar)^2}{2m}\,\frac{\partial^2}{\partial x^2}
  = -\frac{\hbar^2}{2m}\,\frac{\partial^2}{\partial x^2},
  \qquad
  \hat v = \frac{\hat p}{m} = -\frac{\ii\hbar}{m}\,\frac{\partial}{\partial x}.
\]

\begin{center}
\renewcommand{\arraystretch}{1.6}
\begin{tabular}{@{}lll@{}}
\toprule
\textbf{Grandeur} & \textbf{Classique} & \textbf{Opérateur} \\
\midrule
Position & $x$ & $\hat x = x\times$ \\
Quantité de mouvement & $p$ & $\hat p = -\ii\hbar\,\partial_x$ \\
Vitesse & $v=\dfrac{p}{m}$ & $\hat v = -\dfrac{\ii\hbar}{m}\,\partial_x$ \\
Énergie cinétique & $E_c=\dfrac{p^2}{2m}$ & $\hat E_c = -\dfrac{\hbar^2}{2m}\,\partial_x^2$ \\
Énergie potentielle & $V(x)$ & $\hat V = V(x)\times$ \\
Énergie totale & $E=E_c+V$ & $\Hop = \hat E_c+\hat V = -\dfrac{\hbar^2}{2m}\,\partial_x^2+V(x)$ \\
\bottomrule
\end{tabular}
\end{center}

L'opérateur hamiltonien $\Hop$ représente l'énergie totale. Avec $E\to\ii\hbar\,\partial_t$,
l'équation de Schrödinger de la section~\ref{sec:construction} n'est que la relation
$E=E_c+V$ écrite avec des opérateurs : $\ii\hbar\,\partial_t\Psi=\Hop\Psi$.

\begin{remarque}[Convention de signe]
Le signe de $\hat p$ dépend du choix de l'onde plane $\ee^{\ii(kx-\omega t)}$. La
convention conjuguée $\ee^{-\ii(kx-\omega t)}$ inverse les deux signes
($\hat p=+\ii\hbar\,\partial_x$ et $E\to-\ii\hbar\,\partial_t$) sans changer les
résultats physiques. Comme $\hat p$ intervient au carré, $\hat E_c$ est le même dans
les deux conventions.
\end{remarque}

\begin{definition}{Valeur moyenne d'un opérateur}
Pour une fonction d'onde normalisée $\Psi$,
\[
  \langle\hat O\rangle
  = \int_{-\infty}^{+\infty}\Psi^*(x,t)\;\hat O\,\Psi(x,t)\,\dd x
  = \bra{\Psi}\hat O\ket{\Psi}.
\]
C'est la moyenne des résultats obtenus en mesurant la grandeur sur un grand nombre
de systèmes identiques, tous préparés dans l'état $\Psi$. L'opérateur agit sur
$\Psi$ seul, à droite de $\Psi^*$. Pour la position,
$\langle\hat x\rangle=\int x\,|\Psi|^2\,\dd x$ : c'est l'espérance de la densité
de probabilité de Born.

L'\emph{incertitude} (écart type) sur la grandeur associée à $\hat O$ mesure la
dispersion des résultats autour de cette moyenne :
\[
  \sigma_O=\sqrt{\langle\hat O^{\,2}\rangle-\langle\hat O\rangle^2}\,,
  \qquad
  \langle\hat O^{\,2}\rangle=\int\Psi^*\,\hat O^{\,2}\,\Psi\,\dd x .
\]
Si $\Psi$ est une fonction propre de $\hat O$ ($\hat O\Psi=\lambda\Psi$), alors
$\langle\hat O\rangle=\lambda$, $\langle\hat O^{\,2}\rangle=\lambda^2$ et
$\sigma_O=0$ : la grandeur a une valeur certaine.
\end{definition}

% ---------------------------------------------------------------------
\subsection{Solutions de l'équation de Schrödinger}\label{sec:solutions}
% ---------------------------------------------------------------------

\subsubsection{Du potentiel à la fonction d'onde}

En mécanique classique, l'énergie potentielle $V(x,t)$ fixe la force
$F=-\partial V/\partial x$, puis $\dd p/\dd t=F$ donne la trajectoire. En mécanique
quantique, c'est l'équation de Schrödinger qui joue ce rôle : connaissant $V$, on la
résout pour obtenir $\Psi(x,t)$.

\begin{resultat}{ESDT : équation de Schrödinger dépendante du temps}
\[
  \ii\hbar\,\pd{\Psi(x,t)}{t} = \Hop\,\Psi(x,t),
  \qquad
  \Hop = -\frac{\hbar^2}{2m}\,\pdd{}{x} + V(x,t).
\]
\end{resultat}

Si l'énergie potentielle ne dépend pas du temps, $V(x,t)=V(x)$, la force
$F=-\dd V/\dd x$ (dérivée ordinaire) est constante dans le temps et l'énergie se
conserve. L'équation se résout alors par séparation des variables.

\subsubsection{Séparation des variables}

On cherche des solutions produit d'une fonction de $x$ seul et d'une fonction de $t$
seul :
\[
  \Psi(x,t)=\psi(x)\,\varphi(t).
\]
Dans l'ESDT, les dérivées partielles deviennent des dérivées ordinaires :
\[
  \ii\hbar\,\psi(x)\,\frac{\dd\varphi}{\dd t}
  = \varphi(t)\left[-\frac{\hbar^2}{2m}\,\frac{\dd^2\psi}{\dd x^2}+V(x)\,\psi(x)\right].
\]
On divise par $\psi(x)\,\varphi(t)$ :
\[
  \underbrace{\ii\hbar\,\frac{1}{\varphi}\,\frac{\dd\varphi}{\dd t}}_{\text{ne dépend que de }t}
  \;=\;
  \underbrace{\frac{1}{\psi}\left[-\frac{\hbar^2}{2m}\,\frac{\dd^2\psi}{\dd x^2}+V(x)\,\psi\right]}_{\text{ne dépend que de }x}.
\]
Les deux membres dépendent de variables indépendantes et sont pourtant égaux pour
tout $(x,t)$ : ils sont égaux à une même constante, notée $E$. On obtient deux
équations.

\paragraph{Partie temporelle.} $\ii\hbar\,\dd\varphi/\dd t=E\,\varphi$ donne
\[
  \varphi(t)=\ee^{-\ii Et/\hbar}=\ee^{-\ii\omega t},
  \qquad \omega=\frac{E}{\hbar},
\]
la constante multiplicative étant absorbée dans $\psi$. On retrouve $E=\hbar\omega$ :
la constante $E$ est bien l'énergie de la particule.

\paragraph{Partie spatiale.} C'est l'ESIT ci-dessous.

\begin{resultat}{ESIT : équation de Schrödinger indépendante du temps}
\[
  -\frac{\hbar^2}{2m}\,\frac{\dd^2\psi(x)}{\dd x^2}+V(x)\,\psi(x)=E\,\psi(x)
  \qquad\Longleftrightarrow\qquad
  \Hop\,\psi(x)=E\,\psi(x).
\]
C'est un \emph{problème aux valeurs propres} : $E$ est une valeur propre de $\Hop$
(énergie permise) et $\psi$ la fonction propre associée. Chaque solution
$(E_n,\psi_n)$ donne un \emph{état stationnaire}
\[
  \Psi_n(x,t)=\psi_n(x)\,\ee^{-\ii E_nt/\hbar}.
\]
\end{resultat}

Dans un état stationnaire, $|\Psi_n(x,t)|^2=|\psi_n(x)|^2$ ne dépend pas du temps :
le facteur temporel est une phase globale, non observable. Les conditions physiques
($\psi$ continue et normalisable) ne sont satisfaites que pour certaines valeurs de
$E$ : c'est l'origine de la quantification de l'énergie, illustrée ci-dessous.

Comme l'ESDT est linéaire (hypothèse~(c) de la section~\ref{sec:construction}), toute
superposition $\Psi=\sum_n c_n\,\Psi_n$ est aussi solution, normalisée si
$\sum_n|c_n|^2=1$ et si les $\psi_n$ sont orthonormées. Une telle superposition
n'est en général pas stationnaire : $|\Psi|^2$ contient des termes croisés en
$\ee^{\ii(E_m-E_n)t/\hbar}$.

\subsubsection{Exemple : puits de potentiel infini}

Une particule de masse $m$ (un électron par exemple) est enfermée entre deux parois
impénétrables, en $x=0$ et $x=a$ :
\[
  V(x)=
  \begin{cases}
    0 & \text{si } 0<x<a,\\
    +\infty & \text{ailleurs.}
  \end{cases}
\]
La particule ne peut pas se trouver hors du puits : $\psi=0$ pour $x\le0$ et
$x\ge a$. La fonction $\psi$ est continue (mais pas sa dérivée, à cause du saut
infini de $V$), d'où les conditions aux limites
\[
  \psi(0)=\psi(a)=0 .
\]

\paragraph{Opérateur hamiltonien.} Dans le puits $V=0$, donc $\Hop$ se réduit à
l'opérateur énergie cinétique :
\[
  \Hop=\hat E_c=-\frac{\hbar^2}{2m}\,\frac{\dd^2}{\dd x^2}\qquad(0<x<a).
\]

\paragraph{Résolution de l'ESIT.} Avec $V=0$, l'ESIT s'écrit
$\psi''=-k^2\psi$ avec $k=\sqrt{2mE}/\hbar$. On prend $E>0$ : pour $E\le0$, la seule
solution vérifiant $\psi(0)=\psi(a)=0$ est $\psi\equiv0$. La solution générale est
\[
  \psi(x)=A\sin(kx)+B\cos(kx).
\]
\begin{itemize}[nosep]
  \item $\psi(0)=0$ donne $B=0$ ;
  \item $\psi(a)=A\sin(ka)=0$ avec $A\neq0$ (sinon $\psi\equiv0$ : pas de
        particule) donne $ka=n\pi$, soit
        $k_n=\dfrac{n\pi}{a}$ avec $n=1,2,3,\dots$
        ($n=0$ donne $\psi\equiv0$, et $n<0$ redonne les mêmes états au signe près).
\end{itemize}
Cela revient à $\lambda_n=2\pi/k_n=2a/n$, soit $a=n\,\lambda_n/2$ : $n$ demi-longueurs
d'onde tiennent dans le puits, comme pour une onde stationnaire.

\begin{resultat}{Énergie quantifiée}
\[
  E_n=\frac{\hbar^2k_n^2}{2m}=\frac{n^2\pi^2\hbar^2}{2ma^2}=\frac{n^2h^2}{8ma^2}=n^2E_1,
  \qquad n=1,2,3,\dots
\]
\textbf{Énergie fondamentale} (niveau $n=1$, le plus bas) :
$\displaystyle E_1=\frac{\pi^2\hbar^2}{2ma^2}=\frac{h^2}{8ma^2}>0$.
\end{resultat}

\begin{remarque}[Lecture du résultat]
\begin{itemize}[nosep]
  \item Les énergies sont \emph{discrètes} : $E_{n+1}-E_n=(2n+1)E_1$. Comme
        $E_n\propto1/a^2$, plus le puits est étroit, plus les énergies sont grandes.
  \item $E_1\neq0$ : une particule confinée ne peut pas être au repos (énergie de
        point zéro), conséquence du principe d'incertitude de Heisenberg.
  \item Ordre de grandeur : pour un électron avec $a=1$~nm,
        $E_1\approx0{,}38$~eV et $E_2=4E_1\approx1{,}5$~eV. Pour $a=1$~cm,
        $E_1\approx4\times10^{-15}$~eV : la quantification est indétectable.
\end{itemize}
\end{remarque}

\paragraph{Normalisation.} Le module de $\ee^{-\ii E_nt/\hbar}$ vaut 1, donc
$|\Psi_n|^2=A^2\sin^2(n\pi x/a)$ et la condition de Born impose
\[
  \int_0^a|\Psi_n|^2\,\dd x
  = A^2\int_0^a\frac{1-\cos\!\big(2n\pi x/a\big)}{2}\,\dd x
  = A^2\,\frac{a}{2}=1
  \quad\Longrightarrow\quad A=\sqrt{\frac{2}{a}},
\]
car l'intégrale du cosinus sur $n$ périodes est nulle.

\begin{resultat}{États stationnaires du puits infini}
\[
  \Psi_n(x,t)=\sqrt{\frac{2}{a}}\;\sin\!\Big(\frac{n\pi x}{a}\Big)\,\ee^{-\ii E_nt/\hbar}
  \quad(0<x<a),
  \qquad
  |\Psi_n|^2=\frac{2}{a}\,\sin^2\!\Big(\frac{n\pi x}{a}\Big).
\]
$\psi_n$ possède $n-1$ nœuds à l'intérieur du puits. Pour $n=1$ (état fondamental,
d'énergie $E_1$) :
\[
  \Psi_1(x,t)=\sqrt{\frac{2}{a}}\;\sin\!\Big(\frac{\pi x}{a}\Big)\,\ee^{-\ii E_1t/\hbar}.
\]
\end{resultat}

\begin{figure}[!ht]
\centering
\begin{tikzpicture}[>=Stealth, x=1cm, y=0.95cm]
  % ---------- (a) niveaux d'énergie et fonctions propres ----------
  \begin{scope}
    \node[above, font=\small] at (2,5.0) {(a) niveaux $E_n$ et fonctions propres $\psi_n$};
    \draw[->] (-2.5,0) -- (-2.5,4.8) node[above] {$E$};
    \draw[very thick, insarouge] (0,4.4) -- (0,0) -- (4,0) -- (4,4.4);
    \node[insarouge, right, font=\small] at (4.05,3.9) {$V=+\infty$};
    \node[insarouge, below, font=\small] at (2,0) {$V=0$};
    \node[below, font=\small] at (0,-0.35) {$0$};
    \node[below, font=\small] at (4,-0.35) {$a$};
    \node[left, font=\small] at (-0.1,0.4) {$E_1$};
    \node[left, font=\small] at (-0.1,1.6) {$E_2=4E_1$};
    \node[left, font=\small] at (-0.1,3.6) {$E_3=9E_1$};
    \foreach \n/\lev in {1/0.4, 2/1.6, 3/3.6} {
      \draw[dashed, insagris!70] (0,\lev) -- (4,\lev);
      \draw[thick, insagris, domain=0:1, samples=80]
        plot ({4*\x}, {\lev+0.5*sin(\n*180*\x)});
    }
  \end{scope}
  % ---------- (b) densités de probabilité ----------
  \begin{scope}[xshift=7cm]
    \node[above, font=\small] at (2,5.0) {(b) densités de probabilité $|\psi_n|^2$};
    \draw[->] (-0.3,0) -- (4.6,0) node[right] {$x$};
    \draw[thick, insarouge] (0,0) -- (0,4.4) (4,0) -- (4,4.4);
    \node[below, font=\small] at (0,0) {$0$};
    \node[below, font=\small] at (4,0) {$a$};
    \node[below, font=\small] at (2,0) {$a/2$};
    \foreach \n/\lev in {1/0.4, 2/1.6, 3/3.6} {
      \draw[dashed, insagris!70] (0,\lev) -- (4,\lev);
      \fill[insarouge!20] (0,\lev) -- plot[domain=0:1, samples=80]
        ({4*\x}, {\lev+0.6*(sin(\n*180*\x))^2}) -- (4,\lev) -- cycle;
      \draw[thick, insarouge, domain=0:1, samples=80]
        plot ({4*\x}, {\lev+0.6*(sin(\n*180*\x))^2});
    }
    \draw[dotted, thick, insagris] (2,0) -- (2,4.4);
    \node[above, font=\small] at (2,4.4) {$\langle\hat x\rangle=a/2$};
  \end{scope}
\end{tikzpicture}
\caption{Puits de potentiel infini de largeur $a$. (a) Niveaux d'énergie
$E_n=n^2E_1$ et fonctions propres $\psi_n$. (b) Densités de probabilité
$|\psi_n|^2$, symétriques par rapport à $a/2$.}
\end{figure}

\subsubsection{Position moyenne dans le puits}

D'après la définition de la valeur moyenne, avec $\hat x=x\times$ :
\[
  \langle\hat x\rangle
  = \int_0^a\Psi_n^*(x,t)\;x\;\Psi_n(x,t)\,\dd x
  = \frac{2}{a}\int_0^a x\,\sin^2\!\Big(\frac{n\pi x}{a}\Big)\dd x
  = \frac{1}{a}\int_0^a x\Big(1-\cos\frac{2n\pi x}{a}\Big)\dd x .
\]
Le premier terme vaut $\dfrac1a\cdot\dfrac{a^2}{2}=\dfrac a2$. Pour le second, une
intégration par parties donne
\[
  \int_0^a x\cos\frac{2n\pi x}{a}\,\dd x
  =\Big[\frac{a\,x}{2n\pi}\sin\frac{2n\pi x}{a}\Big]_0^a
   +\frac{a^2}{4n^2\pi^2}\Big[\cos\frac{2n\pi x}{a}\Big]_0^a = 0,
\]
puisque $\sin(2n\pi)=0$ et $\cos(2n\pi)=1$.

\begin{resultat}{Position moyenne}
\[
  \langle\hat x\rangle=\frac{a}{2}\qquad\text{pour tout }n\text{ et à tout instant.}
\]
\end{resultat}

Les phases $\ee^{\pm\ii E_nt/\hbar}$ se compensent dans $\Psi^*\Psi$. Le résultat se
voit d'ailleurs sans calcul, car $|\Psi_n|^2$ est symétrique par rapport à $a/2$.

\begin{remarque}[Moyenne et position la plus probable]
Pour $n=2$, la densité de probabilité s'annule en $x=a/2$ alors que
$\langle\hat x\rangle=a/2$ : la valeur moyenne n'est pas la position la plus
probable, et la particule n'est jamais trouvée là.
\end{remarque}

\paragraph{Incertitude sur l'énergie.} Comme $\Hop\Psi_n=E_n\Psi_n$, on a aussi
$\Hop^{2}\Psi_n=E_n^2\Psi_n$, donc $\langle\Hop\rangle=E_n$, $\langle\Hop^{2}\rangle=E_n^2$
et
\[
  \sigma_E=\sqrt{\langle\Hop^{2}\rangle-\langle\Hop\rangle^2}=\sqrt{E_n^2-E_n^2}=0 .
\]
Dans un état propre (stationnaire) l'énergie est parfaitement déterminée, alors que la
position ne l'est pas. Que se passe-t-il si la particule n'est pas dans un état propre ?
C'est l'objet de la section~\ref{sec:superposition}.

% ---------------------------------------------------------------------
\subsection{Bilan : opérateurs et puits infini}
% ---------------------------------------------------------------------

\begin{center}
\small\renewcommand{\arraystretch}{1.25}
\begin{tabular}{@{}ll@{}}
\toprule
\textbf{Notion} & \textbf{Formule} \\
\midrule
Bra-ket & $\braket{a}{b}=a_1^*b_1+a_2^*b_2$, \quad
          $\braket{\Psi}{\Psi}=\displaystyle\int\Psi^*\Psi\,\dd x=1$ \\
Opérateurs & $\hat p=-\ii\hbar\,\partial_x$, \quad
             $\hat E_c=\dfrac{\hat p^{\,2}}{2m}=-\dfrac{\hbar^2}{2m}\,\partial_x^2$, \quad
             $\Hop=\hat E_c+V(x)$ \\
Valeur moyenne, incertitude & $\langle\hat O\rangle=\displaystyle\int\Psi^*\,\hat O\,\Psi\,\dd x$,
          \quad $\sigma_O=\sqrt{\langle\hat O^{\,2}\rangle-\langle\hat O\rangle^2}$ \\
ESDT & $\ii\hbar\,\partial_t\Psi=\Hop\,\Psi$ \\
Séparation ($V=V(x)$) & $\Psi=\psi(x)\,\ee^{-\ii Et/\hbar}$, \quad ESIT : $\Hop\,\psi=E\,\psi$ \\
Puits infini & $\psi_n=\sqrt{2/a}\,\sin(n\pi x/a)$, \quad
               $E_n=\dfrac{n^2\pi^2\hbar^2}{2ma^2}$, \quad $E_1=\dfrac{\pi^2\hbar^2}{2ma^2}$ \\
Moyennes dans $\Psi_n$ & $\langle\hat x\rangle=a/2$, \quad $\langle\Hop\rangle=E_n$, \quad $\sigma_E=0$ \\
\bottomrule
\end{tabular}
\end{center}
.
%  Il PROLONGE chapitre1.tex : pas de \section ici, les sous-sections
%  1.6 à 1.9 font suite à 1.1 à 1.5 (la suite est dans chapitre3.tex).
%  Pas de préambule non plus : les
%  environnements (definition, resultat, remarque), les macros (\ii, \ee,
%  \dd, \pd, \pdd, \Hop, \ket, \bra, \braket) et les couleurs (insarouge,
%  insagris) sont définis dans main.tex.
%  Convention de signe : p = -i hbar d/dx, cohérente avec l'onde plane
%  e^{i(kx-wt)} de la section 1.3 (les notes de cours donnaient +i hbar).
% =====================================================================

% Macros de Dirac (déjà définies dans main.tex récent ; ce bloc évite l'erreur
% « Undefined control sequence \ket » si main.tex est une ancienne version).
\providecommand{\ket}[1]{|#1\rangle}
\providecommand{\bra}[1]{\langle #1|}
\providecommand{\braket}[2]{\langle #1|#2\rangle}

% ---------------------------------------------------------------------
\subsection{Notation de Dirac}\label{sec:dirac}
% ---------------------------------------------------------------------

Un état s'écrit comme un vecteur colonne, par exemple
$\begin{pmatrix} a_1\\ a_2\end{pmatrix}\in\mathbb{C}^2$ : c'est un élément d'un
espace de Hilbert de dimension~2, noté $\mathcal{H}^2$. Dirac le note simplement
$\ket{a}$ et évite d'écrire les composantes à chaque ligne.

\begin{definition}{Notation de Dirac : ket, bra et bra-ket}
\begin{center}
\begin{tabular}{@{}r@{\hspace{0.5em}}l@{\hspace{2em}}l@{}}
  {\Large\color{insarouge}$\ket{a}$}
    & $=\begin{pmatrix}a_1\\ a_2\end{pmatrix}\in\mathcal{H}^2$
    & \textbf{ket} : vecteur colonne \\[1.6ex]
  {\Large\color{insarouge}$\bra{a}$}
    & $=\big(\ket{a}\big)^\dagger=\begin{pmatrix}a_1^* & a_2^*\end{pmatrix}$
    & \textbf{bra} : vecteur ligne \\[1.6ex]
  {\Large\color{insarouge}$\braket{a}{b}$}
    & $=\begin{pmatrix}a_1^* & a_2^*\end{pmatrix}\begin{pmatrix}b_1\\ b_2\end{pmatrix}
       =a_1^*b_1+a_2^*b_2$
    & \textbf{bra-ket} : produit scalaire \\
\end{tabular}
\end{center}
Le bra est le transposé \emph{conjugué} du ket. Le nom vient de l'anglais
\emph{bracket} (crochet) : en accolant $\bra{a}$~(\emph{bra}) et $\ket{b}$~(\emph{ket})
on obtient $\braket{a}{b}$~(\emph{bracket}), un nombre complexe.

Propriétés : $\braket{b}{a}=\braket{a}{b}^*$ et
$\braket{a}{a}=|a_1|^2+|a_2|^2\in\mathbb{R}_+$ (norme au carré).
\end{definition}

Dans l'autre ordre, le produit d'un ket par un bra est une matrice (un opérateur) :
\[
  \ket{a}\bra{b}
  =\begin{pmatrix}a_1\\ a_2\end{pmatrix}\begin{pmatrix}b_1^* & b_2^*\end{pmatrix}
  =\begin{pmatrix}a_1b_1^* & a_1b_2^*\\ a_2b_1^* & a_2b_2^*\end{pmatrix}.
\]

\begin{remarque}[Pourquoi le conjugué ?]
Pour $\ket{a}=\dfrac{1}{\sqrt2}\begin{pmatrix}1\\ \ii\end{pmatrix}$, on a
$\bra{a}=\dfrac{1}{\sqrt2}\begin{pmatrix}1 & -\ii\end{pmatrix}$ et
$\braket{a}{a}=\tfrac12\,(1+1)=1$. Sans conjugaison on trouverait
$\tfrac12\,(1+\ii^2)=0$, ce qui n'est pas une norme.
\end{remarque}

Un état physique est \emph{normalisé} :
\[
  \braket{\Psi}{\Psi}=1 .
\]
Pour une particule dans l'espace, $\ket{\Psi}$ est une fonction de $x$ (dimension
infinie) et le produit scalaire devient une intégrale,
$\braket{\varphi}{\psi}=\int\varphi^*\psi\,\dd x$. On retrouve la condition de
normalisation de Born (section~\ref{sec:born}) :
\[
  \braket{\Psi}{\Psi} = \int_{-\infty}^{+\infty}\Psi^*\Psi\,\dd x = 1 .
\]
L'espace de Hilbert et le formalisme complet sont repris plus loin dans le cours.

% ---------------------------------------------------------------------
\subsection{Opérateurs quantiques}\label{sec:operateurs}
% ---------------------------------------------------------------------

À chaque grandeur physique (position, quantité de mouvement, énergie) on associe
un \emph{opérateur} $\hat O$ : une règle qui transforme une fonction d'onde en une
autre, $\Psi\mapsto\hat O\Psi$. Le point de départ est l'onde plane de la
section~\ref{sec:construction} : $\Psi=\ee^{\ii(kx-\omega t)}$ vérifie
$-\ii\hbar\,\partial_x\Psi=\hbar k\,\Psi=p\,\Psi$.

\begin{definition}{Opérateur quantité de mouvement}
\[
  \hat p = -\ii\hbar\,\frac{\partial}{\partial x},
  \qquad
  \hat p\,\Psi = -\ii\hbar\,\frac{\partial\Psi}{\partial x}.
\]
L'onde plane est fonction propre de $\hat p$, de valeur propre $p=\hbar k$.
\end{definition}

Les autres opérateurs s'obtiennent en remplaçant $p$ par $\hat p$ dans
l'expression classique. Avec $E_c=\tfrac12 mv^2=p^2/2m$ :
\[
  \hat E_c = \frac{\hat p^{\,2}}{2m}
  = \frac{1}{2m}\Big(-\ii\hbar\frac{\partial}{\partial x}\Big)
                \Big(-\ii\hbar\frac{\partial}{\partial x}\Big)
  = \frac{(-\ii\hbar)^2}{2m}\,\frac{\partial^2}{\partial x^2}
  = -\frac{\hbar^2}{2m}\,\frac{\partial^2}{\partial x^2},
  \qquad
  \hat v = \frac{\hat p}{m} = -\frac{\ii\hbar}{m}\,\frac{\partial}{\partial x}.
\]

\begin{center}
\renewcommand{\arraystretch}{1.6}
\begin{tabular}{@{}lll@{}}
\toprule
\textbf{Grandeur} & \textbf{Classique} & \textbf{Opérateur} \\
\midrule
Position & $x$ & $\hat x = x\times$ \\
Quantité de mouvement & $p$ & $\hat p = -\ii\hbar\,\partial_x$ \\
Vitesse & $v=\dfrac{p}{m}$ & $\hat v = -\dfrac{\ii\hbar}{m}\,\partial_x$ \\
Énergie cinétique & $E_c=\dfrac{p^2}{2m}$ & $\hat E_c = -\dfrac{\hbar^2}{2m}\,\partial_x^2$ \\
Énergie potentielle & $V(x)$ & $\hat V = V(x)\times$ \\
Énergie totale & $E=E_c+V$ & $\Hop = \hat E_c+\hat V = -\dfrac{\hbar^2}{2m}\,\partial_x^2+V(x)$ \\
\bottomrule
\end{tabular}
\end{center}

L'opérateur hamiltonien $\Hop$ représente l'énergie totale. Avec $E\to\ii\hbar\,\partial_t$,
l'équation de Schrödinger de la section~\ref{sec:construction} n'est que la relation
$E=E_c+V$ écrite avec des opérateurs : $\ii\hbar\,\partial_t\Psi=\Hop\Psi$.

\begin{remarque}[Convention de signe]
Le signe de $\hat p$ dépend du choix de l'onde plane $\ee^{\ii(kx-\omega t)}$. La
convention conjuguée $\ee^{-\ii(kx-\omega t)}$ inverse les deux signes
($\hat p=+\ii\hbar\,\partial_x$ et $E\to-\ii\hbar\,\partial_t$) sans changer les
résultats physiques. Comme $\hat p$ intervient au carré, $\hat E_c$ est le même dans
les deux conventions.
\end{remarque}

\begin{definition}{Valeur moyenne d'un opérateur}
Pour une fonction d'onde normalisée $\Psi$,
\[
  \langle\hat O\rangle
  = \int_{-\infty}^{+\infty}\Psi^*(x,t)\;\hat O\,\Psi(x,t)\,\dd x
  = \bra{\Psi}\hat O\ket{\Psi}.
\]
C'est la moyenne des résultats obtenus en mesurant la grandeur sur un grand nombre
de systèmes identiques, tous préparés dans l'état $\Psi$. L'opérateur agit sur
$\Psi$ seul, à droite de $\Psi^*$. Pour la position,
$\langle\hat x\rangle=\int x\,|\Psi|^2\,\dd x$ : c'est l'espérance de la densité
de probabilité de Born.

L'\emph{incertitude} (écart type) sur la grandeur associée à $\hat O$ mesure la
dispersion des résultats autour de cette moyenne :
\[
  \sigma_O=\sqrt{\langle\hat O^{\,2}\rangle-\langle\hat O\rangle^2}\,,
  \qquad
  \langle\hat O^{\,2}\rangle=\int\Psi^*\,\hat O^{\,2}\,\Psi\,\dd x .
\]
Si $\Psi$ est une fonction propre de $\hat O$ ($\hat O\Psi=\lambda\Psi$), alors
$\langle\hat O\rangle=\lambda$, $\langle\hat O^{\,2}\rangle=\lambda^2$ et
$\sigma_O=0$ : la grandeur a une valeur certaine.
\end{definition}

% ---------------------------------------------------------------------
\subsection{Solutions de l'équation de Schrödinger}\label{sec:solutions}
% ---------------------------------------------------------------------

\subsubsection{Du potentiel à la fonction d'onde}

En mécanique classique, l'énergie potentielle $V(x,t)$ fixe la force
$F=-\partial V/\partial x$, puis $\dd p/\dd t=F$ donne la trajectoire. En mécanique
quantique, c'est l'équation de Schrödinger qui joue ce rôle : connaissant $V$, on la
résout pour obtenir $\Psi(x,t)$.

\begin{resultat}{ESDT : équation de Schrödinger dépendante du temps}
\[
  \ii\hbar\,\pd{\Psi(x,t)}{t} = \Hop\,\Psi(x,t),
  \qquad
  \Hop = -\frac{\hbar^2}{2m}\,\pdd{}{x} + V(x,t).
\]
\end{resultat}

Si l'énergie potentielle ne dépend pas du temps, $V(x,t)=V(x)$, la force
$F=-\dd V/\dd x$ (dérivée ordinaire) est constante dans le temps et l'énergie se
conserve. L'équation se résout alors par séparation des variables.

\subsubsection{Séparation des variables}

On cherche des solutions produit d'une fonction de $x$ seul et d'une fonction de $t$
seul :
\[
  \Psi(x,t)=\psi(x)\,\varphi(t).
\]
Dans l'ESDT, les dérivées partielles deviennent des dérivées ordinaires :
\[
  \ii\hbar\,\psi(x)\,\frac{\dd\varphi}{\dd t}
  = \varphi(t)\left[-\frac{\hbar^2}{2m}\,\frac{\dd^2\psi}{\dd x^2}+V(x)\,\psi(x)\right].
\]
On divise par $\psi(x)\,\varphi(t)$ :
\[
  \underbrace{\ii\hbar\,\frac{1}{\varphi}\,\frac{\dd\varphi}{\dd t}}_{\text{ne dépend que de }t}
  \;=\;
  \underbrace{\frac{1}{\psi}\left[-\frac{\hbar^2}{2m}\,\frac{\dd^2\psi}{\dd x^2}+V(x)\,\psi\right]}_{\text{ne dépend que de }x}.
\]
Les deux membres dépendent de variables indépendantes et sont pourtant égaux pour
tout $(x,t)$ : ils sont égaux à une même constante, notée $E$. On obtient deux
équations.

\paragraph{Partie temporelle.} $\ii\hbar\,\dd\varphi/\dd t=E\,\varphi$ donne
\[
  \varphi(t)=\ee^{-\ii Et/\hbar}=\ee^{-\ii\omega t},
  \qquad \omega=\frac{E}{\hbar},
\]
la constante multiplicative étant absorbée dans $\psi$. On retrouve $E=\hbar\omega$ :
la constante $E$ est bien l'énergie de la particule.

\paragraph{Partie spatiale.} C'est l'ESIT ci-dessous.

\begin{resultat}{ESIT : équation de Schrödinger indépendante du temps}
\[
  -\frac{\hbar^2}{2m}\,\frac{\dd^2\psi(x)}{\dd x^2}+V(x)\,\psi(x)=E\,\psi(x)
  \qquad\Longleftrightarrow\qquad
  \Hop\,\psi(x)=E\,\psi(x).
\]
C'est un \emph{problème aux valeurs propres} : $E$ est une valeur propre de $\Hop$
(énergie permise) et $\psi$ la fonction propre associée. Chaque solution
$(E_n,\psi_n)$ donne un \emph{état stationnaire}
\[
  \Psi_n(x,t)=\psi_n(x)\,\ee^{-\ii E_nt/\hbar}.
\]
\end{resultat}

Dans un état stationnaire, $|\Psi_n(x,t)|^2=|\psi_n(x)|^2$ ne dépend pas du temps :
le facteur temporel est une phase globale, non observable. Les conditions physiques
($\psi$ continue et normalisable) ne sont satisfaites que pour certaines valeurs de
$E$ : c'est l'origine de la quantification de l'énergie, illustrée ci-dessous.

Comme l'ESDT est linéaire (hypothèse~(c) de la section~\ref{sec:construction}), toute
superposition $\Psi=\sum_n c_n\,\Psi_n$ est aussi solution, normalisée si
$\sum_n|c_n|^2=1$ et si les $\psi_n$ sont orthonormées. Une telle superposition
n'est en général pas stationnaire : $|\Psi|^2$ contient des termes croisés en
$\ee^{\ii(E_m-E_n)t/\hbar}$.

\subsubsection{Exemple : puits de potentiel infini}

Une particule de masse $m$ (un électron par exemple) est enfermée entre deux parois
impénétrables, en $x=0$ et $x=a$ :
\[
  V(x)=
  \begin{cases}
    0 & \text{si } 0<x<a,\\
    +\infty & \text{ailleurs.}
  \end{cases}
\]
La particule ne peut pas se trouver hors du puits : $\psi=0$ pour $x\le0$ et
$x\ge a$. La fonction $\psi$ est continue (mais pas sa dérivée, à cause du saut
infini de $V$), d'où les conditions aux limites
\[
  \psi(0)=\psi(a)=0 .
\]

\paragraph{Opérateur hamiltonien.} Dans le puits $V=0$, donc $\Hop$ se réduit à
l'opérateur énergie cinétique :
\[
  \Hop=\hat E_c=-\frac{\hbar^2}{2m}\,\frac{\dd^2}{\dd x^2}\qquad(0<x<a).
\]

\paragraph{Résolution de l'ESIT.} Avec $V=0$, l'ESIT s'écrit
$\psi''=-k^2\psi$ avec $k=\sqrt{2mE}/\hbar$. On prend $E>0$ : pour $E\le0$, la seule
solution vérifiant $\psi(0)=\psi(a)=0$ est $\psi\equiv0$. La solution générale est
\[
  \psi(x)=A\sin(kx)+B\cos(kx).
\]
\begin{itemize}[nosep]
  \item $\psi(0)=0$ donne $B=0$ ;
  \item $\psi(a)=A\sin(ka)=0$ avec $A\neq0$ (sinon $\psi\equiv0$ : pas de
        particule) donne $ka=n\pi$, soit
        $k_n=\dfrac{n\pi}{a}$ avec $n=1,2,3,\dots$
        ($n=0$ donne $\psi\equiv0$, et $n<0$ redonne les mêmes états au signe près).
\end{itemize}
Cela revient à $\lambda_n=2\pi/k_n=2a/n$, soit $a=n\,\lambda_n/2$ : $n$ demi-longueurs
d'onde tiennent dans le puits, comme pour une onde stationnaire.

\begin{resultat}{Énergie quantifiée}
\[
  E_n=\frac{\hbar^2k_n^2}{2m}=\frac{n^2\pi^2\hbar^2}{2ma^2}=\frac{n^2h^2}{8ma^2}=n^2E_1,
  \qquad n=1,2,3,\dots
\]
\textbf{Énergie fondamentale} (niveau $n=1$, le plus bas) :
$\displaystyle E_1=\frac{\pi^2\hbar^2}{2ma^2}=\frac{h^2}{8ma^2}>0$.
\end{resultat}

\begin{remarque}[Lecture du résultat]
\begin{itemize}[nosep]
  \item Les énergies sont \emph{discrètes} : $E_{n+1}-E_n=(2n+1)E_1$. Comme
        $E_n\propto1/a^2$, plus le puits est étroit, plus les énergies sont grandes.
  \item $E_1\neq0$ : une particule confinée ne peut pas être au repos (énergie de
        point zéro), conséquence du principe d'incertitude de Heisenberg.
  \item Ordre de grandeur : pour un électron avec $a=1$~nm,
        $E_1\approx0{,}38$~eV et $E_2=4E_1\approx1{,}5$~eV. Pour $a=1$~cm,
        $E_1\approx4\times10^{-15}$~eV : la quantification est indétectable.
\end{itemize}
\end{remarque}

\paragraph{Normalisation.} Le module de $\ee^{-\ii E_nt/\hbar}$ vaut 1, donc
$|\Psi_n|^2=A^2\sin^2(n\pi x/a)$ et la condition de Born impose
\[
  \int_0^a|\Psi_n|^2\,\dd x
  = A^2\int_0^a\frac{1-\cos\!\big(2n\pi x/a\big)}{2}\,\dd x
  = A^2\,\frac{a}{2}=1
  \quad\Longrightarrow\quad A=\sqrt{\frac{2}{a}},
\]
car l'intégrale du cosinus sur $n$ périodes est nulle.

\begin{resultat}{États stationnaires du puits infini}
\[
  \Psi_n(x,t)=\sqrt{\frac{2}{a}}\;\sin\!\Big(\frac{n\pi x}{a}\Big)\,\ee^{-\ii E_nt/\hbar}
  \quad(0<x<a),
  \qquad
  |\Psi_n|^2=\frac{2}{a}\,\sin^2\!\Big(\frac{n\pi x}{a}\Big).
\]
$\psi_n$ possède $n-1$ nœuds à l'intérieur du puits. Pour $n=1$ (état fondamental,
d'énergie $E_1$) :
\[
  \Psi_1(x,t)=\sqrt{\frac{2}{a}}\;\sin\!\Big(\frac{\pi x}{a}\Big)\,\ee^{-\ii E_1t/\hbar}.
\]
\end{resultat}

\begin{figure}[!ht]
\centering
\begin{tikzpicture}[>=Stealth, x=1cm, y=0.95cm]
  % ---------- (a) niveaux d'énergie et fonctions propres ----------
  \begin{scope}
    \node[above, font=\small] at (2,5.0) {(a) niveaux $E_n$ et fonctions propres $\psi_n$};
    \draw[->] (-2.5,0) -- (-2.5,4.8) node[above] {$E$};
    \draw[very thick, insarouge] (0,4.4) -- (0,0) -- (4,0) -- (4,4.4);
    \node[insarouge, right, font=\small] at (4.05,3.9) {$V=+\infty$};
    \node[insarouge, below, font=\small] at (2,0) {$V=0$};
    \node[below, font=\small] at (0,-0.35) {$0$};
    \node[below, font=\small] at (4,-0.35) {$a$};
    \node[left, font=\small] at (-0.1,0.4) {$E_1$};
    \node[left, font=\small] at (-0.1,1.6) {$E_2=4E_1$};
    \node[left, font=\small] at (-0.1,3.6) {$E_3=9E_1$};
    \foreach \n/\lev in {1/0.4, 2/1.6, 3/3.6} {
      \draw[dashed, insagris!70] (0,\lev) -- (4,\lev);
      \draw[thick, insagris, domain=0:1, samples=80]
        plot ({4*\x}, {\lev+0.5*sin(\n*180*\x)});
    }
  \end{scope}
  % ---------- (b) densités de probabilité ----------
  \begin{scope}[xshift=7cm]
    \node[above, font=\small] at (2,5.0) {(b) densités de probabilité $|\psi_n|^2$};
    \draw[->] (-0.3,0) -- (4.6,0) node[right] {$x$};
    \draw[thick, insarouge] (0,0) -- (0,4.4) (4,0) -- (4,4.4);
    \node[below, font=\small] at (0,0) {$0$};
    \node[below, font=\small] at (4,0) {$a$};
    \node[below, font=\small] at (2,0) {$a/2$};
    \foreach \n/\lev in {1/0.4, 2/1.6, 3/3.6} {
      \draw[dashed, insagris!70] (0,\lev) -- (4,\lev);
      \fill[insarouge!20] (0,\lev) -- plot[domain=0:1, samples=80]
        ({4*\x}, {\lev+0.6*(sin(\n*180*\x))^2}) -- (4,\lev) -- cycle;
      \draw[thick, insarouge, domain=0:1, samples=80]
        plot ({4*\x}, {\lev+0.6*(sin(\n*180*\x))^2});
    }
    \draw[dotted, thick, insagris] (2,0) -- (2,4.4);
    \node[above, font=\small] at (2,4.4) {$\langle\hat x\rangle=a/2$};
  \end{scope}
\end{tikzpicture}
\caption{Puits de potentiel infini de largeur $a$. (a) Niveaux d'énergie
$E_n=n^2E_1$ et fonctions propres $\psi_n$. (b) Densités de probabilité
$|\psi_n|^2$, symétriques par rapport à $a/2$.}
\end{figure}

\subsubsection{Position moyenne dans le puits}

D'après la définition de la valeur moyenne, avec $\hat x=x\times$ :
\[
  \langle\hat x\rangle
  = \int_0^a\Psi_n^*(x,t)\;x\;\Psi_n(x,t)\,\dd x
  = \frac{2}{a}\int_0^a x\,\sin^2\!\Big(\frac{n\pi x}{a}\Big)\dd x
  = \frac{1}{a}\int_0^a x\Big(1-\cos\frac{2n\pi x}{a}\Big)\dd x .
\]
Le premier terme vaut $\dfrac1a\cdot\dfrac{a^2}{2}=\dfrac a2$. Pour le second, une
intégration par parties donne
\[
  \int_0^a x\cos\frac{2n\pi x}{a}\,\dd x
  =\Big[\frac{a\,x}{2n\pi}\sin\frac{2n\pi x}{a}\Big]_0^a
   +\frac{a^2}{4n^2\pi^2}\Big[\cos\frac{2n\pi x}{a}\Big]_0^a = 0,
\]
puisque $\sin(2n\pi)=0$ et $\cos(2n\pi)=1$.

\begin{resultat}{Position moyenne}
\[
  \langle\hat x\rangle=\frac{a}{2}\qquad\text{pour tout }n\text{ et à tout instant.}
\]
\end{resultat}

Les phases $\ee^{\pm\ii E_nt/\hbar}$ se compensent dans $\Psi^*\Psi$. Le résultat se
voit d'ailleurs sans calcul, car $|\Psi_n|^2$ est symétrique par rapport à $a/2$.

\begin{remarque}[Moyenne et position la plus probable]
Pour $n=2$, la densité de probabilité s'annule en $x=a/2$ alors que
$\langle\hat x\rangle=a/2$ : la valeur moyenne n'est pas la position la plus
probable, et la particule n'est jamais trouvée là.
\end{remarque}

\paragraph{Incertitude sur l'énergie.} Comme $\Hop\Psi_n=E_n\Psi_n$, on a aussi
$\Hop^{2}\Psi_n=E_n^2\Psi_n$, donc $\langle\Hop\rangle=E_n$, $\langle\Hop^{2}\rangle=E_n^2$
et
\[
  \sigma_E=\sqrt{\langle\Hop^{2}\rangle-\langle\Hop\rangle^2}=\sqrt{E_n^2-E_n^2}=0 .
\]
Dans un état propre (stationnaire) l'énergie est parfaitement déterminée, alors que la
position ne l'est pas. Que se passe-t-il si la particule n'est pas dans un état propre ?
C'est l'objet de la section~\ref{sec:superposition}.

% ---------------------------------------------------------------------
\subsection{Bilan : opérateurs et puits infini}
% ---------------------------------------------------------------------

\begin{center}
\small\renewcommand{\arraystretch}{1.25}
\begin{tabular}{@{}ll@{}}
\toprule
\textbf{Notion} & \textbf{Formule} \\
\midrule
Bra-ket & $\braket{a}{b}=a_1^*b_1+a_2^*b_2$, \quad
          $\braket{\Psi}{\Psi}=\displaystyle\int\Psi^*\Psi\,\dd x=1$ \\
Opérateurs & $\hat p=-\ii\hbar\,\partial_x$, \quad
             $\hat E_c=\dfrac{\hat p^{\,2}}{2m}=-\dfrac{\hbar^2}{2m}\,\partial_x^2$, \quad
             $\Hop=\hat E_c+V(x)$ \\
Valeur moyenne, incertitude & $\langle\hat O\rangle=\displaystyle\int\Psi^*\,\hat O\,\Psi\,\dd x$,
          \quad $\sigma_O=\sqrt{\langle\hat O^{\,2}\rangle-\langle\hat O\rangle^2}$ \\
ESDT & $\ii\hbar\,\partial_t\Psi=\Hop\,\Psi$ \\
Séparation ($V=V(x)$) & $\Psi=\psi(x)\,\ee^{-\ii Et/\hbar}$, \quad ESIT : $\Hop\,\psi=E\,\psi$ \\
Puits infini & $\psi_n=\sqrt{2/a}\,\sin(n\pi x/a)$, \quad
               $E_n=\dfrac{n^2\pi^2\hbar^2}{2ma^2}$, \quad $E_1=\dfrac{\pi^2\hbar^2}{2ma^2}$ \\
Moyennes dans $\Psi_n$ & $\langle\hat x\rangle=a/2$, \quad $\langle\Hop\rangle=E_n$, \quad $\sigma_E=0$ \\
\bottomrule
\end{tabular}
\end{center}
.
%  Il PROLONGE chapitre1.tex : pas de \section ici, les sous-sections
%  1.6 à 1.9 font suite à 1.1 à 1.5 (la suite est dans chapitre3.tex).
%  Pas de préambule non plus : les
%  environnements (definition, resultat, remarque), les macros (\ii, \ee,
%  \dd, \pd, \pdd, \Hop, \ket, \bra, \braket) et les couleurs (insarouge,
%  insagris) sont définis dans main.tex.
%  Convention de signe : p = -i hbar d/dx, cohérente avec l'onde plane
%  e^{i(kx-wt)} de la section 1.3 (les notes de cours donnaient +i hbar).
% =====================================================================

% Macros de Dirac (déjà définies dans main.tex récent ; ce bloc évite l'erreur
% « Undefined control sequence \ket » si main.tex est une ancienne version).
\providecommand{\ket}[1]{|#1\rangle}
\providecommand{\bra}[1]{\langle #1|}
\providecommand{\braket}[2]{\langle #1|#2\rangle}

% ---------------------------------------------------------------------
\subsection{Notation de Dirac}\label{sec:dirac}
% ---------------------------------------------------------------------

Un état s'écrit comme un vecteur colonne, par exemple
$\begin{pmatrix} a_1\\ a_2\end{pmatrix}\in\mathbb{C}^2$ : c'est un élément d'un
espace de Hilbert de dimension~2, noté $\mathcal{H}^2$. Dirac le note simplement
$\ket{a}$ et évite d'écrire les composantes à chaque ligne.

\begin{definition}{Notation de Dirac : ket, bra et bra-ket}
\begin{center}
\begin{tabular}{@{}r@{\hspace{0.5em}}l@{\hspace{2em}}l@{}}
  {\Large\color{insarouge}$\ket{a}$}
    & $=\begin{pmatrix}a_1\\ a_2\end{pmatrix}\in\mathcal{H}^2$
    & \textbf{ket} : vecteur colonne \\[1.6ex]
  {\Large\color{insarouge}$\bra{a}$}
    & $=\big(\ket{a}\big)^\dagger=\begin{pmatrix}a_1^* & a_2^*\end{pmatrix}$
    & \textbf{bra} : vecteur ligne \\[1.6ex]
  {\Large\color{insarouge}$\braket{a}{b}$}
    & $=\begin{pmatrix}a_1^* & a_2^*\end{pmatrix}\begin{pmatrix}b_1\\ b_2\end{pmatrix}
       =a_1^*b_1+a_2^*b_2$
    & \textbf{bra-ket} : produit scalaire \\
\end{tabular}
\end{center}
Le bra est le transposé \emph{conjugué} du ket. Le nom vient de l'anglais
\emph{bracket} (crochet) : en accolant $\bra{a}$~(\emph{bra}) et $\ket{b}$~(\emph{ket})
on obtient $\braket{a}{b}$~(\emph{bracket}), un nombre complexe.

Propriétés : $\braket{b}{a}=\braket{a}{b}^*$ et
$\braket{a}{a}=|a_1|^2+|a_2|^2\in\mathbb{R}_+$ (norme au carré).
\end{definition}

Dans l'autre ordre, le produit d'un ket par un bra est une matrice (un opérateur) :
\[
  \ket{a}\bra{b}
  =\begin{pmatrix}a_1\\ a_2\end{pmatrix}\begin{pmatrix}b_1^* & b_2^*\end{pmatrix}
  =\begin{pmatrix}a_1b_1^* & a_1b_2^*\\ a_2b_1^* & a_2b_2^*\end{pmatrix}.
\]

\begin{remarque}[Pourquoi le conjugué ?]
Pour $\ket{a}=\dfrac{1}{\sqrt2}\begin{pmatrix}1\\ \ii\end{pmatrix}$, on a
$\bra{a}=\dfrac{1}{\sqrt2}\begin{pmatrix}1 & -\ii\end{pmatrix}$ et
$\braket{a}{a}=\tfrac12\,(1+1)=1$. Sans conjugaison on trouverait
$\tfrac12\,(1+\ii^2)=0$, ce qui n'est pas une norme.
\end{remarque}

Un état physique est \emph{normalisé} :
\[
  \braket{\Psi}{\Psi}=1 .
\]
Pour une particule dans l'espace, $\ket{\Psi}$ est une fonction de $x$ (dimension
infinie) et le produit scalaire devient une intégrale,
$\braket{\varphi}{\psi}=\int\varphi^*\psi\,\dd x$. On retrouve la condition de
normalisation de Born (section~\ref{sec:born}) :
\[
  \braket{\Psi}{\Psi} = \int_{-\infty}^{+\infty}\Psi^*\Psi\,\dd x = 1 .
\]
L'espace de Hilbert et le formalisme complet sont repris plus loin dans le cours.

% ---------------------------------------------------------------------
\subsection{Opérateurs quantiques}\label{sec:operateurs}
% ---------------------------------------------------------------------

À chaque grandeur physique (position, quantité de mouvement, énergie) on associe
un \emph{opérateur} $\hat O$ : une règle qui transforme une fonction d'onde en une
autre, $\Psi\mapsto\hat O\Psi$. Le point de départ est l'onde plane de la
section~\ref{sec:construction} : $\Psi=\ee^{\ii(kx-\omega t)}$ vérifie
$-\ii\hbar\,\partial_x\Psi=\hbar k\,\Psi=p\,\Psi$.

\begin{definition}{Opérateur quantité de mouvement}
\[
  \hat p = -\ii\hbar\,\frac{\partial}{\partial x},
  \qquad
  \hat p\,\Psi = -\ii\hbar\,\frac{\partial\Psi}{\partial x}.
\]
L'onde plane est fonction propre de $\hat p$, de valeur propre $p=\hbar k$.
\end{definition}

Les autres opérateurs s'obtiennent en remplaçant $p$ par $\hat p$ dans
l'expression classique. Avec $E_c=\tfrac12 mv^2=p^2/2m$ :
\[
  \hat E_c = \frac{\hat p^{\,2}}{2m}
  = \frac{1}{2m}\Big(-\ii\hbar\frac{\partial}{\partial x}\Big)
                \Big(-\ii\hbar\frac{\partial}{\partial x}\Big)
  = \frac{(-\ii\hbar)^2}{2m}\,\frac{\partial^2}{\partial x^2}
  = -\frac{\hbar^2}{2m}\,\frac{\partial^2}{\partial x^2},
  \qquad
  \hat v = \frac{\hat p}{m} = -\frac{\ii\hbar}{m}\,\frac{\partial}{\partial x}.
\]

\begin{center}
\renewcommand{\arraystretch}{1.6}
\begin{tabular}{@{}lll@{}}
\toprule
\textbf{Grandeur} & \textbf{Classique} & \textbf{Opérateur} \\
\midrule
Position & $x$ & $\hat x = x\times$ \\
Quantité de mouvement & $p$ & $\hat p = -\ii\hbar\,\partial_x$ \\
Vitesse & $v=\dfrac{p}{m}$ & $\hat v = -\dfrac{\ii\hbar}{m}\,\partial_x$ \\
Énergie cinétique & $E_c=\dfrac{p^2}{2m}$ & $\hat E_c = -\dfrac{\hbar^2}{2m}\,\partial_x^2$ \\
Énergie potentielle & $V(x)$ & $\hat V = V(x)\times$ \\
Énergie totale & $E=E_c+V$ & $\Hop = \hat E_c+\hat V = -\dfrac{\hbar^2}{2m}\,\partial_x^2+V(x)$ \\
\bottomrule
\end{tabular}
\end{center}

L'opérateur hamiltonien $\Hop$ représente l'énergie totale. Avec $E\to\ii\hbar\,\partial_t$,
l'équation de Schrödinger de la section~\ref{sec:construction} n'est que la relation
$E=E_c+V$ écrite avec des opérateurs : $\ii\hbar\,\partial_t\Psi=\Hop\Psi$.

\begin{remarque}[Convention de signe]
Le signe de $\hat p$ dépend du choix de l'onde plane $\ee^{\ii(kx-\omega t)}$. La
convention conjuguée $\ee^{-\ii(kx-\omega t)}$ inverse les deux signes
($\hat p=+\ii\hbar\,\partial_x$ et $E\to-\ii\hbar\,\partial_t$) sans changer les
résultats physiques. Comme $\hat p$ intervient au carré, $\hat E_c$ est le même dans
les deux conventions.
\end{remarque}

\begin{definition}{Valeur moyenne d'un opérateur}
Pour une fonction d'onde normalisée $\Psi$,
\[
  \langle\hat O\rangle
  = \int_{-\infty}^{+\infty}\Psi^*(x,t)\;\hat O\,\Psi(x,t)\,\dd x
  = \bra{\Psi}\hat O\ket{\Psi}.
\]
C'est la moyenne des résultats obtenus en mesurant la grandeur sur un grand nombre
de systèmes identiques, tous préparés dans l'état $\Psi$. L'opérateur agit sur
$\Psi$ seul, à droite de $\Psi^*$. Pour la position,
$\langle\hat x\rangle=\int x\,|\Psi|^2\,\dd x$ : c'est l'espérance de la densité
de probabilité de Born.

L'\emph{incertitude} (écart type) sur la grandeur associée à $\hat O$ mesure la
dispersion des résultats autour de cette moyenne :
\[
  \sigma_O=\sqrt{\langle\hat O^{\,2}\rangle-\langle\hat O\rangle^2}\,,
  \qquad
  \langle\hat O^{\,2}\rangle=\int\Psi^*\,\hat O^{\,2}\,\Psi\,\dd x .
\]
Si $\Psi$ est une fonction propre de $\hat O$ ($\hat O\Psi=\lambda\Psi$), alors
$\langle\hat O\rangle=\lambda$, $\langle\hat O^{\,2}\rangle=\lambda^2$ et
$\sigma_O=0$ : la grandeur a une valeur certaine.
\end{definition}

% ---------------------------------------------------------------------
\subsection{Solutions de l'équation de Schrödinger}\label{sec:solutions}
% ---------------------------------------------------------------------

\subsubsection{Du potentiel à la fonction d'onde}

En mécanique classique, l'énergie potentielle $V(x,t)$ fixe la force
$F=-\partial V/\partial x$, puis $\dd p/\dd t=F$ donne la trajectoire. En mécanique
quantique, c'est l'équation de Schrödinger qui joue ce rôle : connaissant $V$, on la
résout pour obtenir $\Psi(x,t)$.

\begin{resultat}{ESDT : équation de Schrödinger dépendante du temps}
\[
  \ii\hbar\,\pd{\Psi(x,t)}{t} = \Hop\,\Psi(x,t),
  \qquad
  \Hop = -\frac{\hbar^2}{2m}\,\pdd{}{x} + V(x,t).
\]
\end{resultat}

Si l'énergie potentielle ne dépend pas du temps, $V(x,t)=V(x)$, la force
$F=-\dd V/\dd x$ (dérivée ordinaire) est constante dans le temps et l'énergie se
conserve. L'équation se résout alors par séparation des variables.

\subsubsection{Séparation des variables}

On cherche des solutions produit d'une fonction de $x$ seul et d'une fonction de $t$
seul :
\[
  \Psi(x,t)=\psi(x)\,\varphi(t).
\]
Dans l'ESDT, les dérivées partielles deviennent des dérivées ordinaires :
\[
  \ii\hbar\,\psi(x)\,\frac{\dd\varphi}{\dd t}
  = \varphi(t)\left[-\frac{\hbar^2}{2m}\,\frac{\dd^2\psi}{\dd x^2}+V(x)\,\psi(x)\right].
\]
On divise par $\psi(x)\,\varphi(t)$ :
\[
  \underbrace{\ii\hbar\,\frac{1}{\varphi}\,\frac{\dd\varphi}{\dd t}}_{\text{ne dépend que de }t}
  \;=\;
  \underbrace{\frac{1}{\psi}\left[-\frac{\hbar^2}{2m}\,\frac{\dd^2\psi}{\dd x^2}+V(x)\,\psi\right]}_{\text{ne dépend que de }x}.
\]
Les deux membres dépendent de variables indépendantes et sont pourtant égaux pour
tout $(x,t)$ : ils sont égaux à une même constante, notée $E$. On obtient deux
équations.

\paragraph{Partie temporelle.} $\ii\hbar\,\dd\varphi/\dd t=E\,\varphi$ donne
\[
  \varphi(t)=\ee^{-\ii Et/\hbar}=\ee^{-\ii\omega t},
  \qquad \omega=\frac{E}{\hbar},
\]
la constante multiplicative étant absorbée dans $\psi$. On retrouve $E=\hbar\omega$ :
la constante $E$ est bien l'énergie de la particule.

\paragraph{Partie spatiale.} C'est l'ESIT ci-dessous.

\begin{resultat}{ESIT : équation de Schrödinger indépendante du temps}
\[
  -\frac{\hbar^2}{2m}\,\frac{\dd^2\psi(x)}{\dd x^2}+V(x)\,\psi(x)=E\,\psi(x)
  \qquad\Longleftrightarrow\qquad
  \Hop\,\psi(x)=E\,\psi(x).
\]
C'est un \emph{problème aux valeurs propres} : $E$ est une valeur propre de $\Hop$
(énergie permise) et $\psi$ la fonction propre associée. Chaque solution
$(E_n,\psi_n)$ donne un \emph{état stationnaire}
\[
  \Psi_n(x,t)=\psi_n(x)\,\ee^{-\ii E_nt/\hbar}.
\]
\end{resultat}

Dans un état stationnaire, $|\Psi_n(x,t)|^2=|\psi_n(x)|^2$ ne dépend pas du temps :
le facteur temporel est une phase globale, non observable. Les conditions physiques
($\psi$ continue et normalisable) ne sont satisfaites que pour certaines valeurs de
$E$ : c'est l'origine de la quantification de l'énergie, illustrée ci-dessous.

Comme l'ESDT est linéaire (hypothèse~(c) de la section~\ref{sec:construction}), toute
superposition $\Psi=\sum_n c_n\,\Psi_n$ est aussi solution, normalisée si
$\sum_n|c_n|^2=1$ et si les $\psi_n$ sont orthonormées. Une telle superposition
n'est en général pas stationnaire : $|\Psi|^2$ contient des termes croisés en
$\ee^{\ii(E_m-E_n)t/\hbar}$.

\subsubsection{Exemple : puits de potentiel infini}

Une particule de masse $m$ (un électron par exemple) est enfermée entre deux parois
impénétrables, en $x=0$ et $x=a$ :
\[
  V(x)=
  \begin{cases}
    0 & \text{si } 0<x<a,\\
    +\infty & \text{ailleurs.}
  \end{cases}
\]
La particule ne peut pas se trouver hors du puits : $\psi=0$ pour $x\le0$ et
$x\ge a$. La fonction $\psi$ est continue (mais pas sa dérivée, à cause du saut
infini de $V$), d'où les conditions aux limites
\[
  \psi(0)=\psi(a)=0 .
\]

\paragraph{Opérateur hamiltonien.} Dans le puits $V=0$, donc $\Hop$ se réduit à
l'opérateur énergie cinétique :
\[
  \Hop=\hat E_c=-\frac{\hbar^2}{2m}\,\frac{\dd^2}{\dd x^2}\qquad(0<x<a).
\]

\paragraph{Résolution de l'ESIT.} Avec $V=0$, l'ESIT s'écrit
$\psi''=-k^2\psi$ avec $k=\sqrt{2mE}/\hbar$. On prend $E>0$ : pour $E\le0$, la seule
solution vérifiant $\psi(0)=\psi(a)=0$ est $\psi\equiv0$. La solution générale est
\[
  \psi(x)=A\sin(kx)+B\cos(kx).
\]
\begin{itemize}[nosep]
  \item $\psi(0)=0$ donne $B=0$ ;
  \item $\psi(a)=A\sin(ka)=0$ avec $A\neq0$ (sinon $\psi\equiv0$ : pas de
        particule) donne $ka=n\pi$, soit
        $k_n=\dfrac{n\pi}{a}$ avec $n=1,2,3,\dots$
        ($n=0$ donne $\psi\equiv0$, et $n<0$ redonne les mêmes états au signe près).
\end{itemize}
Cela revient à $\lambda_n=2\pi/k_n=2a/n$, soit $a=n\,\lambda_n/2$ : $n$ demi-longueurs
d'onde tiennent dans le puits, comme pour une onde stationnaire.

\begin{resultat}{Énergie quantifiée}
\[
  E_n=\frac{\hbar^2k_n^2}{2m}=\frac{n^2\pi^2\hbar^2}{2ma^2}=\frac{n^2h^2}{8ma^2}=n^2E_1,
  \qquad n=1,2,3,\dots
\]
\textbf{Énergie fondamentale} (niveau $n=1$, le plus bas) :
$\displaystyle E_1=\frac{\pi^2\hbar^2}{2ma^2}=\frac{h^2}{8ma^2}>0$.
\end{resultat}

\begin{remarque}[Lecture du résultat]
\begin{itemize}[nosep]
  \item Les énergies sont \emph{discrètes} : $E_{n+1}-E_n=(2n+1)E_1$. Comme
        $E_n\propto1/a^2$, plus le puits est étroit, plus les énergies sont grandes.
  \item $E_1\neq0$ : une particule confinée ne peut pas être au repos (énergie de
        point zéro), conséquence du principe d'incertitude de Heisenberg.
  \item Ordre de grandeur : pour un électron avec $a=1$~nm,
        $E_1\approx0{,}38$~eV et $E_2=4E_1\approx1{,}5$~eV. Pour $a=1$~cm,
        $E_1\approx4\times10^{-15}$~eV : la quantification est indétectable.
\end{itemize}
\end{remarque}

\paragraph{Normalisation.} Le module de $\ee^{-\ii E_nt/\hbar}$ vaut 1, donc
$|\Psi_n|^2=A^2\sin^2(n\pi x/a)$ et la condition de Born impose
\[
  \int_0^a|\Psi_n|^2\,\dd x
  = A^2\int_0^a\frac{1-\cos\!\big(2n\pi x/a\big)}{2}\,\dd x
  = A^2\,\frac{a}{2}=1
  \quad\Longrightarrow\quad A=\sqrt{\frac{2}{a}},
\]
car l'intégrale du cosinus sur $n$ périodes est nulle.

\begin{resultat}{États stationnaires du puits infini}
\[
  \Psi_n(x,t)=\sqrt{\frac{2}{a}}\;\sin\!\Big(\frac{n\pi x}{a}\Big)\,\ee^{-\ii E_nt/\hbar}
  \quad(0<x<a),
  \qquad
  |\Psi_n|^2=\frac{2}{a}\,\sin^2\!\Big(\frac{n\pi x}{a}\Big).
\]
$\psi_n$ possède $n-1$ nœuds à l'intérieur du puits. Pour $n=1$ (état fondamental,
d'énergie $E_1$) :
\[
  \Psi_1(x,t)=\sqrt{\frac{2}{a}}\;\sin\!\Big(\frac{\pi x}{a}\Big)\,\ee^{-\ii E_1t/\hbar}.
\]
\end{resultat}

\begin{figure}[!ht]
\centering
\begin{tikzpicture}[>=Stealth, x=1cm, y=0.95cm]
  % ---------- (a) niveaux d'énergie et fonctions propres ----------
  \begin{scope}
    \node[above, font=\small] at (2,5.0) {(a) niveaux $E_n$ et fonctions propres $\psi_n$};
    \draw[->] (-2.5,0) -- (-2.5,4.8) node[above] {$E$};
    \draw[very thick, insarouge] (0,4.4) -- (0,0) -- (4,0) -- (4,4.4);
    \node[insarouge, right, font=\small] at (4.05,3.9) {$V=+\infty$};
    \node[insarouge, below, font=\small] at (2,0) {$V=0$};
    \node[below, font=\small] at (0,-0.35) {$0$};
    \node[below, font=\small] at (4,-0.35) {$a$};
    \node[left, font=\small] at (-0.1,0.4) {$E_1$};
    \node[left, font=\small] at (-0.1,1.6) {$E_2=4E_1$};
    \node[left, font=\small] at (-0.1,3.6) {$E_3=9E_1$};
    \foreach \n/\lev in {1/0.4, 2/1.6, 3/3.6} {
      \draw[dashed, insagris!70] (0,\lev) -- (4,\lev);
      \draw[thick, insagris, domain=0:1, samples=80]
        plot ({4*\x}, {\lev+0.5*sin(\n*180*\x)});
    }
  \end{scope}
  % ---------- (b) densités de probabilité ----------
  \begin{scope}[xshift=7cm]
    \node[above, font=\small] at (2,5.0) {(b) densités de probabilité $|\psi_n|^2$};
    \draw[->] (-0.3,0) -- (4.6,0) node[right] {$x$};
    \draw[thick, insarouge] (0,0) -- (0,4.4) (4,0) -- (4,4.4);
    \node[below, font=\small] at (0,0) {$0$};
    \node[below, font=\small] at (4,0) {$a$};
    \node[below, font=\small] at (2,0) {$a/2$};
    \foreach \n/\lev in {1/0.4, 2/1.6, 3/3.6} {
      \draw[dashed, insagris!70] (0,\lev) -- (4,\lev);
      \fill[insarouge!20] (0,\lev) -- plot[domain=0:1, samples=80]
        ({4*\x}, {\lev+0.6*(sin(\n*180*\x))^2}) -- (4,\lev) -- cycle;
      \draw[thick, insarouge, domain=0:1, samples=80]
        plot ({4*\x}, {\lev+0.6*(sin(\n*180*\x))^2});
    }
    \draw[dotted, thick, insagris] (2,0) -- (2,4.4);
    \node[above, font=\small] at (2,4.4) {$\langle\hat x\rangle=a/2$};
  \end{scope}
\end{tikzpicture}
\caption{Puits de potentiel infini de largeur $a$. (a) Niveaux d'énergie
$E_n=n^2E_1$ et fonctions propres $\psi_n$. (b) Densités de probabilité
$|\psi_n|^2$, symétriques par rapport à $a/2$.}
\end{figure}

\subsubsection{Position moyenne dans le puits}

D'après la définition de la valeur moyenne, avec $\hat x=x\times$ :
\[
  \langle\hat x\rangle
  = \int_0^a\Psi_n^*(x,t)\;x\;\Psi_n(x,t)\,\dd x
  = \frac{2}{a}\int_0^a x\,\sin^2\!\Big(\frac{n\pi x}{a}\Big)\dd x
  = \frac{1}{a}\int_0^a x\Big(1-\cos\frac{2n\pi x}{a}\Big)\dd x .
\]
Le premier terme vaut $\dfrac1a\cdot\dfrac{a^2}{2}=\dfrac a2$. Pour le second, une
intégration par parties donne
\[
  \int_0^a x\cos\frac{2n\pi x}{a}\,\dd x
  =\Big[\frac{a\,x}{2n\pi}\sin\frac{2n\pi x}{a}\Big]_0^a
   +\frac{a^2}{4n^2\pi^2}\Big[\cos\frac{2n\pi x}{a}\Big]_0^a = 0,
\]
puisque $\sin(2n\pi)=0$ et $\cos(2n\pi)=1$.

\begin{resultat}{Position moyenne}
\[
  \langle\hat x\rangle=\frac{a}{2}\qquad\text{pour tout }n\text{ et à tout instant.}
\]
\end{resultat}

Les phases $\ee^{\pm\ii E_nt/\hbar}$ se compensent dans $\Psi^*\Psi$. Le résultat se
voit d'ailleurs sans calcul, car $|\Psi_n|^2$ est symétrique par rapport à $a/2$.

\begin{remarque}[Moyenne et position la plus probable]
Pour $n=2$, la densité de probabilité s'annule en $x=a/2$ alors que
$\langle\hat x\rangle=a/2$ : la valeur moyenne n'est pas la position la plus
probable, et la particule n'est jamais trouvée là.
\end{remarque}

\paragraph{Incertitude sur l'énergie.} Comme $\Hop\Psi_n=E_n\Psi_n$, on a aussi
$\Hop^{2}\Psi_n=E_n^2\Psi_n$, donc $\langle\Hop\rangle=E_n$, $\langle\Hop^{2}\rangle=E_n^2$
et
\[
  \sigma_E=\sqrt{\langle\Hop^{2}\rangle-\langle\Hop\rangle^2}=\sqrt{E_n^2-E_n^2}=0 .
\]
Dans un état propre (stationnaire) l'énergie est parfaitement déterminée, alors que la
position ne l'est pas. Que se passe-t-il si la particule n'est pas dans un état propre ?
C'est l'objet de la section~\ref{sec:superposition}.

% ---------------------------------------------------------------------
\subsection{Bilan : opérateurs et puits infini}
% ---------------------------------------------------------------------

\begin{center}
\small\renewcommand{\arraystretch}{1.25}
\begin{tabular}{@{}ll@{}}
\toprule
\textbf{Notion} & \textbf{Formule} \\
\midrule
Bra-ket & $\braket{a}{b}=a_1^*b_1+a_2^*b_2$, \quad
          $\braket{\Psi}{\Psi}=\displaystyle\int\Psi^*\Psi\,\dd x=1$ \\
Opérateurs & $\hat p=-\ii\hbar\,\partial_x$, \quad
             $\hat E_c=\dfrac{\hat p^{\,2}}{2m}=-\dfrac{\hbar^2}{2m}\,\partial_x^2$, \quad
             $\Hop=\hat E_c+V(x)$ \\
Valeur moyenne, incertitude & $\langle\hat O\rangle=\displaystyle\int\Psi^*\,\hat O\,\Psi\,\dd x$,
          \quad $\sigma_O=\sqrt{\langle\hat O^{\,2}\rangle-\langle\hat O\rangle^2}$ \\
ESDT & $\ii\hbar\,\partial_t\Psi=\Hop\,\Psi$ \\
Séparation ($V=V(x)$) & $\Psi=\psi(x)\,\ee^{-\ii Et/\hbar}$, \quad ESIT : $\Hop\,\psi=E\,\psi$ \\
Puits infini & $\psi_n=\sqrt{2/a}\,\sin(n\pi x/a)$, \quad
               $E_n=\dfrac{n^2\pi^2\hbar^2}{2ma^2}$, \quad $E_1=\dfrac{\pi^2\hbar^2}{2ma^2}$ \\
Moyennes dans $\Psi_n$ & $\langle\hat x\rangle=a/2$, \quad $\langle\Hop\rangle=E_n$, \quad $\sigma_E=0$ \\
\bottomrule
\end{tabular}
\end{center}
